%% !TEX root = ../manual_gaolv5.tex
% Chapter 12.
%% Chapter 12: How GAOL computes its intervals. Sources: gaol.tex of f5bf180 lines
%% 1142-1150 and 1369-1423; doc/simd.md (whole); doc/compare/README.md and
%% performance.md (whole); the introductions and summaries of
%% doc/compare/special_cases.md, enclosure.md, itf1788.md and tang2021.md.
\chapter{How GAOL computes its intervals}\label{chap:computing}

The chapters before this one describe what GAOL computes: the operations,
their bounds and their accuracy. This chapter describes how GAOL computes
them, and what each part of the computation costs. An interval library
spends its time in three places: the arithmetic itself, which is one or a
few instructions; the rounding direction, which the processor holds in a
control register that is slow to write; and the elementary functions. Which
of the three weighs most depends on the processor, on the way GAOL is built
and on the program. The measurements of this chapter say how much each
weighs on the machines we had, and they are given in the manner of an
experimental study: each table of measurements names its machine and, where the source gives them, its compiler and its date,
and no number is given that the measurements do not give.

\medskip
In this chapter, we first present the representation of an interval in a
register and the kernels that compute with it on each family of processors,
SSE2, AVX-512, NEON and the floating-point unit, with the rule that makes a
build take the AVX-512 path (Section~\ref{sec:kernels}).
Section~\ref{sec:where-time-goes} profiles the demonstration programs of GAOL
to see where the time of a program goes.
Section~\ref{sec:rounding-by-processor} gives the cost of a write of the
control register on each processor and what each instruction set offers
instead, and Section~\ref{sec:compilers-give} what the compilers give, and
do not give, to a program that needs a rounding direction.
Section~\ref{sec:comparison} compares GAOL v5 with GAOL 4.2.3 and with five
other libraries of interval arithmetic, on the special cases of the standard,
on the enclosure of the elementary functions and on the time of the
operations. Section~\ref{sec:test-suites} reports the run on GAOL v5 of ITF1788,
a suite of test cases for IEEE 1788-2015 written outside the standard, and of
the benchmark of Tang et al., and Section~\ref{sec:new-paths} draws the
verdicts on the paths that were considered and says what a new one has to
show. The chapter closes with a discussion.

%%--------------------------------------------------------------------------
\section{The representation of an interval and the kernels}
\label{sec:kernels}
%%--------------------------------------------------------------------------
\cindex{kernel}\cindex{SSE2}\cindex{NEON}\cindex{AVX-512}%
\cindex{register of an interval}%

On most platforms, switching the rounding direction is costly. The number of
switches is cut down by Proposition~\ref{prop:round-neg}
(\pxref{sec:containment}): GAOL stores the lower bound of an interval negated,
and computes both bounds rounding upward, so that the rounding direction can be
set toward $+\infty$ once and for all.

\begin{Def}\label{def:computing-register}
An interval $\itv{a}{b}$ is kept as the pair \code{<-a, b>}, its lower
bound negated, in one vector of two doubles: a \code{\_\_m128d} on x86 with
the SSE2 intervals, a \code{float64x2\_t} on 64-bit ARM with the NEON
intervals. The FPU intervals keep the same two numbers in two doubles. The
\dfn{kernel} of an operation is the function of the library that computes
the two doubles of its result from the two doubles of each operand, which it
takes and gives in registers, after checking that the rounding direction is
upward and setting it when it is not (the kernels of the AVX-512 path carry
the direction in the instruction, set none, and check only the modes that flush the subnormals to zero, \pxref{sec:avx512}).
\end{Def}

\begin{Ex}
Let $x = \itv{a}{b}$ and $y = \itv{c}{d}$ be kept as \code{<-a, b>} and
\code{<-c, d>}. One vector addition rounded upward gives
$\langle \roundUp{-a-c}, \roundUp{b+d} \rangle$, which is
$\langle -\roundDn{a+c}, \roundUp{b+d} \rangle$ by
Proposition~\ref{prop:round-neg}: the register of $x+y$. The lower
bound is minus the upper bound of the opposite, so that one addition rounded
upward gives both bounds (the accuracy of each operation is the subject of
Chapter~\ref{chap:accuracy}).
\end{Ex}

\medskip
Table~\ref{tab:computing-kernels} says what each build computes with vector
instructions and what stays scalar. The SSE2 intervals are those of x86, the
NEON intervals those of 64-bit ARM, both under the option \option{GAOL\_SIMD}
of CMake (\pxref{sec:neon}); the AVX-512 path is a variant of the SSE2
intervals (\pxref{sec:avx512}); the FPU intervals are those of every other
processor (\pxref{sec:fpu-intervals}).

\begin{table}[htbp]
\centering\footnotesize
\caption[What each build of GAOL computes with vector instructions]{What each build of GAOL computes with vector instructions, and what stays scalar (from \file{doc/simd.md}).} \label{tab:computing-kernels}
\begin{tabular}{|L{3.6cm}|L{4.8cm}|L{4.5cm}|}
\hline
\textbf{Build} & \textbf{Vector kernels} & \textbf{Scalar} \\
\hline
SSE2 (x86, \code{GAOL\_SIMD}) &
\code{+ - * / sqr}, the operations with a double, \code{inverse},
\code{\%=}, \code{-x}, the integer powers with \code{\_\_FMA\_\_}, the
leaves of \code{div\_rel} &
\code{sqrt}, \code{\&=}, \code{|=}, \code{mid}, \code{midpoint}, the
relations, the elementary functions, the reverse functions \\
\hline
AVX-512 path &
\code{+ - * / sqr}, the operations with a double and the integer powers
(with \code{\_\_FMA\_\_}), with the rounding direction in the instruction;
\code{sqrt} as two scalar square roots and \code{/} as two scalar divisions,
with embedded rounding &
as SSE2 but \code{sqrt}; \code{inverse}, \code{\%=} and the leaves of \code{div\_rel} stay the vector kernels of SSE2, which set the direction \\
\hline
NEON (arm64, \code{GAOL\_SIMD}; every compiler) &
\code{+ - * / sqr sqrt}, the operations with a double, \code{\%=},
\code{inverse}, \code{\&=}, \code{|=}, \code{-x}, the integer powers (two
lanes with \code{vfmaq\_f64}) &
\code{mid}, \code{midpoint}, \code{div\_rel}, the relations, the elementary functions, the reverse functions, and \code{sqrt} of an upper bound above $2^{1000}$ \\
\hline
FPU (the rest) & none & everything \\
\hline
\end{tabular}

\end{table}

\begin{Rem}
The two bounds of an interval are one vector of two doubles, and \code{+},
\code{-}, \code{*}, \code{/}, \code{sqr} (and \code{sqrt} where it has a
kernel) are one or a few vector instructions: a wider register does not help
one interval. Only a batch interface, on boxes or arrays of intervals, which
GAOL has not, could use 256 or 512 bits.
\end{Rem}

\subsection{The SSE2 and NEON kernels}

The NEON intervals compute as the SSE2 intervals do on
x86.\newinvfive{} The sum and the difference are one vector addition rounded
upward. The product is the maximum of the four products of the bounds, with
no branch on their signs. The quotient is one vector division when the
divisor has no zero, and single divisions otherwise. With NEON, the square roots of
both bounds are one vector square root rounded upward, the lower bound being
that root where it is exact, and the double below it otherwise. The compilers
that take the NEON intervals, the options that select them and what the user
sees of them are described in Chapter~\ref{chap:rounding} (\pxref{sec:neon}).

\medskip
The products and the quotients were nine cases of the signs of the bounds, and the squares three, told apart by a \code{switch} in GAOL 4.2.3 and in GAOL v5 at
first. When the signs of the operands are random, as those of the benchmark
of Section~\ref{sec:comparison} are, a mispredicted branch costs more than
the operation it chooses. The kernels now compute with no branch on the
signs:\newinvfive{} the larger of four products for a product, one division
of the two lanes by the denominators the signs choose, by masks, for a
quotient, two squares and a mask for a square. The two doubles of each operand
are passed in a register to the function of the library that computes an
operation, which returns the two doubles of its result in one, where the
operators of GAOL 4.2.3 read their operand and wrote their result through
memory, the result going back to memory for the operation that followed.
Each operation still checks that the rounding direction is upward, with one
comparison, and sets it when it is not, where GAOL 4.2.3, in its trust mode, takes it upward on
trust (\pxref{sec:rounding-direction}).

\medskip
The NEON intervals check the direction by reading the control register FPCR
with \code{mrs}, which raises no exception flag, and write it with
\code{msr} only if the direction is not upward with the flush modes clear
(\pxref{sec:ftz}); the FPU intervals add $1+(2^{-1060}+0)$ instead, which
raises the inexact flag (\pxref{sec:fp-exceptions}). \code{msr} on FPCR is a
serializing instruction, slower than \code{mrs}, which is why it is written
when needed only. With \option{GAOL\_PRESERVE\_ROUNDING}, each operation
reads FPCR once and writes back the direction and the modes it found at its
end only if they differ from what the operation left: a program that
computes upward already, with the flush modes clear, pays no write
(\pxref{sec:preserve-rounding}). The bounds of the NEON intervals are those
of the FPU intervals but for the sign of a zero bound (\pxref{sec:zero-sign}),
which the maximum of 64-bit ARM, ranking $+0$ above $-0$, can set otherwise
than the SSE2 and the FPU intervals do, where the \code{maxpd} of x86
returns its second operand. The tool \file{tests/tools/arith} compares the
NEON and the FPU intervals on a grid of operands, with the calling code in
each of the four rounding directions: it found no other difference.

\subsection{When the build takes the AVX-512 path}
\label{sec:computing-avx512-rule}
\findex{GAOL\_PREFER\_AVX512}\findex{GAOL\_PRESERVE\_ROUNDING}%

The AVX-512 path is off in the default build. The build takes it only where
the option \option{GAOL\_PRESERVE\_ROUNDING} is on, where the compiler has the
instructions and where the machine building runs them, and leaves it out
otherwise; the setting of the developers \option{-DGAOL\_PREFER\_AVX512=ON}
asks for it without that option, and \option{=OFF} refuses it
(\pxref{sec:avx512}). A library built with the path needs a processor with the
instructions. The
build decides it the same way with the three build
systems.\newinvfive{} For \code{+}, \code{-}, \code{*} and \code{sqr} the
path puts the two bounds in two of the eight lanes of a 512-bit register,
for the embedded rounding only; \code{/} and \code{sqrt} are scalar
instructions with the embedded rounding, two of them for the two bounds
(\pxref{sec:new-paths}).

\medskip
The rule comes from four builds of one source, \option{GAOL\_PRESERVE\_ROUNDING}
on or off, each with the path (\option{GAOL\_PREFER\_AVX512} on) and without
it, measured by the script \cmd{performance-ab.sh} of \file{.github/scripts/}
with \option{--runs 9 --lines}: \cmd{gaol\_performance} and the long lines of
\file{doc/compare/code/run\_lines.sh}, the two builds run in turn, the median
of 9 runs, in nanoseconds, on an Intel i7-1185G7, whose AVX-512 units are
halved, with Clang 21.1.8, \option{-O3 -mfma}, on 7 November 2026, the
machine doing nothing else. Table~\ref{tab:computing-avx512} gives the
times; the ratio is the time with SSE2 over the time with the path, so that
above 1 the path is faster.

\begin{table}[htbp]
\centering\footnotesize
\caption[The SSE2 intervals against the AVX-512 path]{The SSE2 intervals against the AVX-512 path, with and without \option{GAOL\_PRESERVE\_ROUNDING}, in nanoseconds per operation or per line (median of 9 runs, Intel i7-1185G7, Clang 21.1.8, 7 November 2026).} \label{tab:computing-avx512}
\begin{tabular}{|l|r|r|r|r|r|r|}
\hline
& \multicolumn{3}{c|}{\textbf{Not preserved}}
& \multicolumn{3}{c|}{\textbf{Preserved}} \\
\textbf{Operation} & SSE2 & AVX-512 & Ratio & SSE2 & AVX-512 & Ratio \\
\hline
\code{x + y} & 1.22 & 1.51 & 0.81 & 8.82 & 1.49 & 5.92 \\
\code{x - y} & 1.33 & 1.56 & 0.85 & 9.38 & 1.53 & 6.13 \\
\code{x * y} & 2.75 & 4.45 & 0.62 & 12.14 & 4.46 & 2.72 \\
\code{x / y} & 3.21 & 5.00 & 0.64 & 11.05 & 5.07 & 2.18 \\
\code{sqr(x)} & 1.97 & 2.03 & 0.97 & 8.42 & 2.04 & 4.13 \\
\code{sqrt(x)} & 8.52 & 4.71 & 1.81 & 13.54 & 4.72 & 2.87 \\
\code{pow(x, 3)} & 6.61 & 7.93 & 0.83 & 10.42 & 7.83 & 1.33 \\
\hline
line of 15 \code{+} (\code{sum16}) & 19.67 & 40.66 & 0.48 & 176.0 & 37.92
& 4.64 \\
line of 7 \code{*} (\code{mul8}) & 23.80 & 37.42 & 0.64 & 86.56 & 35.14
& 2.46 \\
line of 7 \code{/} (\code{quot8}) & 26.50 & 46.56 & 0.57 & 86.73 & 43.59
& 1.99 \\
line of 13 \code{+ - * /} (\code{mix13}) & 27.75 & 43.23 & 0.64 & 152.8
& 41.49 & 3.68 \\
\hline
\end{tabular}

\end{table}

The median of the seven operations is 0.83 without
\option{GAOL\_PRESERVE\_ROUNDING}: the operations take 20\,\% longer with
the path, and the ten long lines 72\,\% longer (the table gives four of them;
the others are in the output of the script); their median ratio is 0.58; only
\code{sqrt} gains. It is 2.87 with it: every operation gains, the ten lines by
1.91 to 4.64, with a median of 3.19. The path pays where each operation would otherwise write the
control register twice, and costs where the register is written only if the
direction is not upward already, as it is by default: the rule is therefore
that of \option{GAOL\_PRESERVE\_ROUNDING}. The threshold was a median gain of
5\,\% over \code{+ - * / sqr sqrt} and \code{pown}. One processor was
measured, and not one with full-width AVX-512 units: such a processor is
expected to lose less without \option{GAOL\_PRESERVE\_ROUNDING}, which the
rule does not rely on.

\subsection{CORE-MATH is scalar}
\cindex{CORE-MATH!scalar}\cindex{fused multiply-add}\findex{GAOL\_FMA}%

CORE-MATH \cite{CoreMath:2022}, in \file{3rd/math-core}, is scalar, with no
vector variant; GAOL calls 38 of its binary64 functions, once per bound.
The fused multiply-add, \option{GAOL\_FMA}, pays in these functions and in
the integer powers of GAOL, \code{cos} being 2.5 times faster and
\code{pow(x, 3)} 3.2 times, and does nothing for the arithmetic; the cost of
each build option is given with the installation
(\pxref{sec:install-options-worth}).

%%--------------------------------------------------------------------------
\section{Where the time of a program goes}
\label{sec:where-time-goes}
%%--------------------------------------------------------------------------
\cindex{profile of a program}%

To know which part of GAOL a program waits for, the demonstration programs of
\file{examples/} were profiled. The call counts of
Table~\ref{tab:computing-profile} are exact, from a NEON build with
\option{-finstrument-functions}. The shares of the time are indicative to
about 3 points: one Apple M2 (NEON intervals), the profiler \cmd{sample} of
macOS at 1 ms during 3 s per program, each sample put to the nearest GAOL
function above it, other work running on the machine; the two passes of the
profile differ by up to 3 points, and the table gives the second.

\begin{table}[htbp]
\centering\footnotesize
\caption[Where the time of the demonstration programs goes]{Where the time of the demonstration programs goes: the calls of GAOL and the shares of the time in the arithmetic, in the elementary functions and outside GAOL (Apple M2, NEON intervals, the profiler \cmd{sample} of macOS; from \file{doc/simd.md}).} \label{tab:computing-profile}
\begin{tabular}{|L{2.9cm}|L{2.6cm}|L{2.6cm}|r|r|r|}
\hline
\textbf{Program} & \textbf{Calls} (\code{+ - * /}, \code{sqr}, powers)
& \textbf{Elementary calls} & \textbf{Arith.} & \textbf{Elem.}
& \textbf{Outside} \\
\hline
03 dependency problem & 0.14 M, 26 k & none & 73\,\% & 0 & 16\,\% \\
04 centered form & 10 k & cos 10 k, acos 8 k & 9\,\% & 82\,\% & 6\,\% \\
06 global optimization & 25 M, 5 M, 0.49 M & none & 31\,\% & 0 & 63\,\% \\
08 SIVIA & 1.0 M, 0.5 M & sin 173 k & 14\,\% & 12\,\% & 71\,\% \\
09 estimation & 0.5 M, 68 k & exp 247 k & 20\,\% & 42\,\% & 31\,\% \\
11 ODE & 0.2 M, 8 k, 4.9 k & sin 12 k, cos 3 k, exp 1 k & 29\,\% & 34\,\%
& 30\,\% \\
16 Goldstein-Price & 32 M, 6 M powers & none & 98\,\% & 0 & 2\,\% \\
\hline
\end{tabular}

\end{table}

\begin{itemize}
\item[$\bullet$] \textbf{The solvers.} On the six solvers 06 to 11, the
arithmetic is 28\,\% of the time weighted by the duration of a run (20\,\%
with equal weights), the elementary functions 4\,\% (20\,\%), and the code
outside GAOL (the priority queue, \code{malloc} of the boxes, the program
itself) 62\,\% (45\,\%). The weighted figures are those of 06, which lasts
80\,\% of the total. Inside GAOL, at equal weights, about 40\,\% is
arithmetic, 34\,\% the elementary functions, 10\,\% the reading and writing
of text and 7\,\% \code{mid()}.
\item[$\bullet$] \textbf{The code around CORE-MATH.} The code of GAOL around
CORE-MATH weighs 21 to 66\,\% of the time of an elementary function (median
about 42\,\%): the check of the rounding direction, the reduction of the
argument, the construction of the result. In the code of \code{exp} around
the CORE-MATH call, 43\,\% of the samples fall on the line of the check of
the direction, a \code{mrs} of 0.8 ns: the samples probably land on the
instruction that waits for it. The CORE-MATH calls themselves cost 3 to
13 ns on the M2.
\item[$\bullet$] \textbf{The other operations.} The relations, \code{hull},
\code{intersect}, \code{min} and \code{max} are 3.6\,\% or less in every
program; \code{mid()} is 6\,\% of the global optimization, and
\code{midpoint()} 11 to 12\,\% of the Krawczyk iteration (29 ns on the
M2: two writes of FPCR).
\item[$\bullet$] \textbf{The check of the direction.} The check (one
\code{mrs} of FPCR here), removed in a copy, saves 0.8\,\% (03), 1.4\,\%
(09), 4.7\,\% (06) and 8.3\,\% (16). \option{GAOL\_PRESERVE\_ROUNDING} on the
M2, for a program that computes to nearest, makes 03 and 16 about 16 times
slower, 06 eight times, 09 five times. On x86 with the AVX-512 path the
ratio is not measured for these programs; the operations of the benchmark of Section~\ref{sec:computing-performance} take 0.6 to 1.5 times as long (Table~\ref{tab:computing-times}).
\item[$\bullet$] \textbf{The text.} The text is 41 to 69\,\% of the five
demonstration programs that print (01, 02, 05, 12, 15; 28 to 31\,\% of 14 and
11\,\% of 03), at 0.2 to 0.45 $\mu$s per bound written. Among the solvers it
is 0 to 1\,\% of 06, 08, 09 and 11, 12\,\% of 10 and 22\,\% of 07.
\end{itemize}

\medskip
\noindent On an Intel i7-1185G7, in the benchmark of
Section~\ref{sec:computing-performance} (a million operations on arrays of
intervals, so partly bound by the memory), the time of an operation is 2.2 ns
for \code{+} and \code{-}, 2.30 for \code{*}, 2.52 for \code{/}, 1.87 for
\code{sqr}, 6.30 for \code{sqrt} and 5.12 for \code{pow(x, 3)}, the same
operation on doubles taking 0.96 ns for \code{+}. With
\option{GAOL\_PRESERVE\_ROUNDING} and the AVX-512 path: 2.19, 2.16, 3.42,
3.60, 2.04, 3.75 and 6.18. On the M2, measured with \file{tests/performance.cpp} (the median of nine runs, Table~\ref{tab:install-worth-arm}), in which the i7-1185G7 takes 1.07~ns for \code{x + y} (the best of nine runs,
Table~\ref{tab:install-worth-x86}; 1.22~ns, the median of nine runs, in
Table~\ref{tab:computing-avx512}), the NEON intervals take 0.90 ns for
\code{+}, 1.52 for \code{*}, 2.08 for \code{/}, 1.27 for \code{sqr} and 2.04
for \code{sqrt}, against 1.15, 3.26, 3.22, 1.78 and 5.17 with the FPU
intervals, whose \code{*}, \code{/} and \code{sqr} take no branch on the
signs (1.15, 3.25, 3.98, 2.07 and 4.49 before).

%%--------------------------------------------------------------------------
\section{The rounding direction, by processor}
\label{sec:rounding-by-processor}
%%--------------------------------------------------------------------------
\cindex{rounding direction!cost of a write}\cindex{MXCSR}\cindex{FPCR}%

Except on the AVX-512 path, an operation computes upward after a write of
the control register of the processor when the direction of the program is
not upward already. What a write costs, by family:

\begin{itemize}
\item[$\bullet$] \textbf{x86.} \code{LDMXCSR} takes 1 to 6 cycles of
throughput on the performance cores of Intel and on AMD processors (1.0 on
Zen 5, 6.0 on an Arrow Lake P-core) and 8 on the efficiency cores, according
to uops.info, which measures one instruction at a time and so does not count
the wait for the operations in flight. \code{STMXCSR}, the read, takes 1
cycle on most Intel cores but 14 to 16 on Zen 2 to Zen 5: one more reason,
with those given in \file{gaol/gaol\_fpu.h} (800 ns under Rosetta 2, 80 ns
with 32-bit Visual C++, a single read reused by Clang 18), why the check of
GAOL on x86-64, in the default build, is an addition and not a read. Two writes took 6 ns a
\code{sqr} (8.3 ns, 4.4 times, in the benchmark of
Section~\ref{sec:computing-performance}) on the SSE2 path with \option{GAOL\_PRESERVE\_ROUNDING}, which the AVX-512 path spares.
\item[$\bullet$] \textbf{64-bit ARM.} \code{mrs} of FPCR takes 0.8 ns,
\code{msr} 4 ns for the value FPCR holds and 11 ns when it changes the
direction (Apple M2). The reference manuals of the Neoverse N1, N2, N3 and
of the Cortex-A78 say that a write that changes the control fields
introduces a barrier or a flush.
\item[$\bullet$] \textbf{32-bit ARM, POWER, s390x, LoongArch, MIPS, RISC-V.}
No figure was found in the references.
\end{itemize}

\medskip
\noindent Table~\ref{tab:computing-processors} says what each instruction
set offers instead of a write, read from its reference manual (the versions
are listed at the end of \file{doc/simd.md}): whether an instruction can
carry its own rounding direction, whether it has a fused multiply-add, which
register holds the direction, and what a path of GAOL could be.

\begingroup\footnotesize
\begin{longtable}{|L{2.0cm}|L{3.8cm}|L{2.1cm}|L{2.0cm}|L{2.9cm}|}
\caption[The rounding direction on each processor family]{The rounding direction on each processor family: what the
instruction set offers, and what a path of GAOL could be (from
\file{doc/simd.md}).}\label{tab:computing-processors}\\
\hline
\textbf{Processor} & \textbf{Direction carried by the instruction}
& \textbf{FMA} & \textbf{Register} & \textbf{What a path could be} \\
\hline
\endfirsthead
\multicolumn{5}{l}{\textit{Table~\thetable, continued}}\\*
\hline
\textbf{Processor} & \textbf{Direction carried by the instruction}
& \textbf{FMA} & \textbf{Register} & \textbf{What a path could be} \\
\hline
\endhead
\hline
\endfoot
\endlastfoot
x86-64 SSE2 & no & FMA3, a set of its own: AVX2 does not imply it
& MXCSR (RC, FTZ, DAZ) & the SSE2 path of today \\
\hline
x86-64 AVX-512F & yes, \code{\{er\}}, for the scalar instructions and the
512-bit ones & yes & MXCSR: RC and the exception flags are ignored, FTZ and
DAZ are honoured & the AVX-512 path of today \\
\hline
x86-64 AVX-512VL, AVX10.1, AVX10.2 & no at 128 and 256 bits:
\code{VADDPD ymm} has no \code{\{er\}} in the Intel manual; AVX10.2 had it
for 256 bits in revisions 1 to 3 and removed it in revision 4.0 (8 May
2025), and revision 7.0 (6 April 2026) keeps it for 512 bits only & yes &
MXCSR & none; a vector of two doubles has to go through the 512-bit form,
as the AVX-512 path does \\
\hline
AArch64, NEON, SVE, SME & no: only \code{FCVT*}, \code{FRINT*} and
\code{FCVTX*} name a rounding in the A64 instruction set of March 2026;
\code{FADD}, \code{FMUL}, \code{FDIV}, \code{FSQRT}, \code{FMLA} use FPCR,
scalar, Advanced SIMD and SVE & yes & FPCR (RMode, FZ; FEAT\_AFP adds FIZ,
AH and NEP) & the NEON path of today; SVE has 128-bit vectors on the
Neoverse N2, N3, V2, V3 and the Cortex-A720, 256 on the V1, and adds no
rounding \\
\hline
ARMv7 & no; NEON has no double, and always rounds to nearest & VFPv4 only
& FPSCR (RMode, FZ) & the FPU intervals \\
\hline
POWER (VSX) & no: \code{xsadddp} and \code{xvadddp} use FPSCR.RN; the
\code{RO} bit exists for the quad precision only (round to odd) & yes &
FPSCR (RN, NI) & VSX \code{<-l, r>} with \code{mffscrni}, needs hardware to
measure \\
\hline
s390x & no: \code{ADB} and \code{VFA} have no rounding field; the \code{M3}
field of the conversions is not the arithmetic & yes & FPC (BRM) &
\code{VFADB}, needs hardware to measure \\
\hline
RISC-V, scalar (D) & yes: the \code{rm} field of \code{FADD.D},
\code{FMUL.D}, \code{FDIV.D}, \code{FSQRT.D}, \code{FMADD.D}\ldots;
\code{rup} is 011 & yes & frm (no flush-to-zero mode) & two scalar
instructions, no check, no register to restore; inline assembly only \\
\hline
RISC-V, vector (V) & no: all the vector floating-point operations use frm;
the \code{\_rm} intrinsics compile to \code{fsrmi} and \code{fsrm} around
the operation & yes & frm & none \\
\hline
LoongArch & no & yes & FCSR & none \\
\hline
MIPS & no (read from the Linux headers and the compiler output, not from the
manual) & not verified & FCSR (RM, FS) & none \\
\hline
WebAssembly 3.0 & no: ``Non-default directed rounding attributes are not
supported'' & \code{relaxed\_madd} may not be fused & none & none: only the
rounding to nearest \\
\hline
\end{longtable}
\endgroup

The RISC-V scalar path is the one case of a rounding direction by
instruction that GAOL has not. Its interest is that it needs no check and no
restore, so that \option{GAOL\_PRESERVE\_ROUNDING} would cost nothing, as on
the AVX-512 path. No figure exists for it: no RISC-V board is available. The
boards read about, in secondary sources, are the SpacemiT K1, which has the
vector extension 1.0, the SiFive P550, which has none, and the T-Head C910,
which has the draft 0.7.1.

%%--------------------------------------------------------------------------
\section{What the compilers give}
\label{sec:compilers-give}
%%--------------------------------------------------------------------------
\cindex{FENV\_ROUND@\texttt{FENV\_ROUND}}%
\cindex{-frounding-math@\option{-frounding-math}}%

No compiler gives a rounding direction by instruction from C or C++: the
paths of GAOL are written with intrinsics or inline assembly. What the
compilers offer was tried, or read from their output, as follows.

\begin{itemize}
\item[$\bullet$] \textbf{The pragma \code{FENV\_ROUND}.}
\code{\#pragma STDC FENV\_ROUND} (C23) would ask for a static rounding. On
the six targets tried (x86-64, aarch64, riscv64, ppc64le, s390x,
loongarch64), GCC 14 to 16 and its trunk ignore it, with a warning only
under \option{-Wall} (GCC bug 20785, open since 2005); Clang 18 to its trunk
parses it, warns ``not supported'' and emits the ordinary instruction (the
IR carries \code{round.upward}; no back end of the three tried, x86-64,
aarch64 and riscv64, uses it); Visual C++ warns C4068. It does not exist in
C++. \option{-Werror=unknown-pragmas} makes Clang's ``not supported'' an
error, but not its other forms of the same pragma (``expected identifier''
with the \code{FE\_UPWARD} of \file{<fenv.h>}, which is a number; ``not
supported on this target''), which belong to another group of warnings, and
\option{/we4068} should do it for Visual C++: none of these was tried.
\item[$\bullet$] \textbf{C++.} P2746R7 (Boehm, January 2025) would replace
\code{fesetround} by a type with an explicit mode, without wording. P4231R0
(May 2026, target C++29) unifies it with P3864; no revision since. Nothing
is in C++26, and none proposes a SIMD form.
\item[$\bullet$] \textbf{The intrinsics of the embedded rounding.}
\code{\_mm\_add\_round\_sd} and \code{\_mm512\_add\_round\_pd} (and the
other operations) compile with GCC 14 to its trunk, Clang 18 to its trunk
and Visual C++ 19.20 to 19.51. The mode must be a literal, and the compilers
do not fold a constant computed with \code{\{ru-sae\}}. There is no
\code{\_mm\_add\_round\_pd}, and no \code{\_mm256\_add\_round\_pd} with GCC
15 and 16 or Clang 21 and later; Clang 20 with \option{-mavx10.2} declares
it and compiles it to the encoding of the withdrawn revisions of AVX10.2,
which the current Intel manual makes \code{\#UD} (never run: no processor
with AVX10.2 is on sale, as far as we found). In inline assembly, the GNU
\cmd{as} of GCC 16.2 on Compiler Explorer assembles
\code{vaddpd \{ru-sae\}, \%ymm1, \%ymm0, \%ymm2} (bytes
\code{62 f1 f9 58 58 d1}, against \code{62 f1 fd 58 58 d1} for the 512-bit
form): no assembly of that form may be written.
\item[$\bullet$] \textbf{\option{-frounding-math}.} GCC does not honour it
across a change of the direction (bugs 34678 and 123482, open): on
\code{fesetround(UP); r = a + b; fesetround(old)}, GCC 14 to its trunk moves
the addition after the restore, on x86-64 and on AArch64. Clang keeps it
between. The vector intrinsics are ordinary floating-point operations for
them: with Apple clang 21 on the M2, at \option{-O2} with
\option{-frounding-math} or \option{-ffp-model=strict}, \code{vaddq\_f64}
and \code{vfmaq\_f64} are right, \code{vpaddq\_f64}, \code{vaddvq\_f64} and
\code{vmulxq\_f64} are computed once, in the first direction, and at
\option{-O2} alone everything is (GAOL uses none of these three). What held
on every compiler tried (GCC 14 to its trunk and Clang 18 to its trunk,
compiled; Apple clang 21, run on the M2) is an empty assembly statement on
each operand after the write of the control register, and on the result, or
a copy of it through volatile memory, before the next one. GAOL does that:
\code{reread\_bounds()} and \code{rnd\_reread()} are empty \code{asm}
statements, \code{GAOL\_RND\_KEEP()} is a copy through volatile memory with
SSE2 and the FPU and an empty \code{asm} with NEON, and
\code{GAOL\_RND\_BARRIER()} is a memory barrier. A path written with
intrinsics needs the same. Visual C++ for arm64 has no \code{asm}
statement: its NEON intervals copy the result through volatile memory,
\code{GAOL\_RND\_BARRIER()} and the writes of FPCR are
\code{\_ReadWriteBarrier()}, and the operands rely on \option{/fp:strict}
(checked by the generated code and by \file{tests/rounding\_direction.cpp},
not by the language).
\item[$\bullet$] \textbf{The intrinsics of the other families.} NEON
\code{float64x2\_t} with GCC, Clang and Visual C++ for arm64 from 19.26; SVE
\file{<arm\_sve.h>} with GCC 14 and Clang 18, and in Visual C++ 19.50 and
19.51 only partially (\code{svmla\_f64\_x} becomes an unresolved call
without a diagnostic); VSX \code{vec\_add} with \option{-mvsx}; RVV with
GCC from 14.1 and Clang from 17.0.1. \code{std::experimental::simd} is in
libstdc++ (libc++ has the types but no \code{operator+}), and
\code{std::simd} of C++26 only in GCC 16, for x86-64.
\end{itemize}

\begin{Rem}
A program that relies on \code{\#pragma STDC FENV\_ROUND} compiles and
rounds in the wrong direction, without an error; so do the
\code{\_mm256\_add\_round\_pd} of Clang 20 and the 256-bit \code{\{ru-sae\}}
assembly. A path of GAOL therefore trusts no rounding direction that the
generated code and \file{tests/rounding\_direction.cpp} have not checked.
\end{Rem}

%%--------------------------------------------------------------------------
\section{GAOL compared with the other libraries}
\label{sec:comparison}
%%--------------------------------------------------------------------------
\cindex{libieeep1788}\cindex{filib++}\cindex{PROFIL/BIAS}%
\cindex{Solaris Studio}\cindex{Boost.Interval}\cindex{comparison of libraries}%

Six implementations of interval arithmetic on doubles are compared in
\file{doc/compare/}: libieeep1788, GAOL v5, filib++, the intervals of
Solaris Studio, PROFIL/BIAS and Boost.Interval.
Table~\ref{tab:computing-libraries} says what each is, in which language it
is written, how it computes its bounds and which arithmetic it follows.
GAOL v5 gives its intervals under two namespaces, \code{gaol}, its own
names, compatible with GAOL~4, and \code{gaol\_ieee1788}, the names of IEEE
1788-2015 \cite{IEEE1788:2015}. They compute with the same functions, and
differ mainly in the power: the \code{pow} of \code{gaol} takes the integer
power for an integer exponent, a negative base included, where the
\code{pow} of \code{gaol\_ieee1788} is the real power of the standard,
defined for $x > 0$ only, and \code{pown} the integer power
(\pxref{sec:pow-namespaces}). The directory \file{doc/compare/code/} holds
the scripts and programs, in C++, Fortran, Python and bash, that run the
three parts of the comparison again and rewrite the tables of the reports.

\begingroup\footnotesize
\begin{longtable}{|L{2.5cm}|L{1.9cm}|L{4.7cm}|L{4.0cm}|}
\caption[The six implementations compared]{The six implementations compared: what they are, their language,
how their bounds are computed and the arithmetic they follow (from
\file{doc/compare/README.md}).}\label{tab:computing-libraries}\\
\hline
\textbf{Library} & \textbf{Language} & \textbf{Bounds}
& \textbf{Arithmetic} \\
\hline
\endfirsthead
\multicolumn{4}{l}{\textit{Table~\thetable, continued}}\\*
\hline
\textbf{Library} & \textbf{Language} & \textbf{Bounds}
& \textbf{Arithmetic} \\
\hline
\endhead
\hline
\endfoot
\endlastfoot
libieeep1788, by Marco Nehmeier (University of W\"urzburg), last commit in
2015 & C++11, header-only & Computed by MPFR, correctly rounded & The
set-based flavor of the preliminary IEEE P1788, the prototype of IEEE
1788-2015 \\
\hline
GAOL v5, by Jordan Ninin (ENSTA), which continues the GAOL of Fr\'ed\'eric
Goualard & C++ & Computed with the rounding direction set upward; elementary
functions with CORE-MATH, correctly rounded upward & Set-based, following
IEEE 1788-2015 in most of its special cases, with the deviations the reports
list \\
\hline
filib++ 3.0.2.2 (University of Wuppertal), as IBEX distributes it, which
IBEX computes with by default on Windows & C++, templates and a small
library & Computed with the rounding direction set and restored by each
operation (\code{native\_switched}); elementary functions of its own & The
extended mode of filib++ (\code{i\_mode\_extended\_flag}, as IBEX uses it),
where the empty set and the infinities are handled \\
\hline
Solaris Studio: the \code{interval(8)} type of Sun's Fortran 95 compiler,
\cmd{f90 -xia}, in Solaris Studio 12.4 (2014) & Fortran 95, intervals built
into the compiler & Computed by \code{libsunimath} & The containment sets of
Sun's interval arithmetic (G.~W. Walster), where the infinities are
values \\
\hline
PROFIL/BIAS 2.0.8 (2009), by Olaf Kn\"uppel and Christian Keil (TU
Hamburg-Harburg) & C++ (PROFIL) over C (BIAS) & Computed by the BIAS
routines, which set the rounding direction downward, then upward, then back
to nearest; elementary functions from the libm, moved outward & The interval
arithmetic before IEEE 1788: no empty set, and an argument outside the
domain of a function is an error that aborts the program \\
\hline
Boost.Interval of Boost 1.92.0 (2026), by Herv\'e Br\"onnimann, Guillaume
Melquiond and Sylvain Pion, the interval library C++ users find first &
C++03 templates, header-only & Computed with the rounding direction set
upward by each operation, the lower bounds negated, and set back after it
(the default policies,
\code{save\_state<}\allowbreak\code{rounded\_arith\_opp>}); elementary
functions from the C library, called with the rounding direction set
downward or upward (\code{rounded\_transc\_opp}), a direction glibc does not
always honour & The interval arithmetic before IEEE 1788, without empty set:
its default checking (\code{checking\_strict}) throws an exception where an
operation would create one \\
\hline
\end{longtable}
\endgroup

The comparison has three parts, each with a report of its own in
\file{doc/compare/}:
\begin{itemize}
\item[$\bullet$] \textbf{Special cases.} 290 special cases taken from
GAOL's tests (infinities, zeros, NaN, empty sets, divisions by zero,
\code{pow} and \code{pown}, \code{+=} and the other operators with doubles,
reading from text, midpoints, widths, radii, comparisons and \code{atan2}),
computed by the six libraries and compared with the results of IEEE
1788-2015, GAOL under its own names and under those of the standard
(Section~\ref{sec:computing-special-cases}).
\item[$\bullet$] \textbf{Enclosure.} Whether the elementary functions of
GAOL and Boost.Interval enclose their values, at 100\,000 random arguments
for each function, against mpmath, GAOL under both namespaces
(Section~\ref{sec:computing-enclosure}).
\item[$\bullet$] \textbf{Performance.} The time of a million additions,
subtractions, multiplications, divisions, sines, cosines, powers, one-line
combinations, evaluations of Shekel 5 and five-line blocks, and the time of
long lines of $+$, $-$, $\times$ and $\div$ between 8 intervals, GAOL v5
being called under the names of the standard, as the other libraries of IEEE
1788, and GAOL 4.2.3 under its own
(Section~\ref{sec:computing-performance}).
\end{itemize}

\medskip
\noindent Table~\ref{tab:computing-inshort} sums the three parts up. Its
times were measured on 17 October 2026, every library in the same run,
compiled by Clang 21.1, on the machine of
Section~\ref{sec:computing-performance}.

\begin{table}[htbp]
\centering\footnotesize
\setlength{\tabcolsep}{3pt}
\caption[The comparison in short]{The comparison in short: the special cases, the enclosure of the elementary functions and the time of six operations (17 October 2026, every library in the same run, Clang 21.1, Intel i7-1185G7).} \label{tab:computing-inshort}
\begin{tabular}{|L{2.15cm}|L{1.5cm}|L{2.2cm}|L{1.3cm}|L{1.6cm}|L{1.7cm}|
  L{2.0cm}|}
\hline
& \textbf{libieee\allowbreak p1788} & \textbf{GAOL} & \textbf{filib++}
& \textbf{Solaris Studio} & \textbf{PROFIL/\allowbreak BIAS}
& \textbf{Boost.\allowbreak Interval} \\
\hline
Special cases with IEEE 1788's result & 281 of 281 & 288 of 288 under
\code{gaol\_ieee1788}; 280 under \code{gaol} & 126 of 245 & 150 of 261
& 56 of 223 & 93 of 168 \\
\ldots\ or an interval enclosing it & 0 & 0; 2 & 50 & 41 & 22 & 9 \\
\ldots\ or another result & 0 & 0; 6 & 69 & 70 & 145 & 66 \\
Cases it has no operation for & 8 & 0 & 44 & 28 & 65 & 121 \\
\hline
Elementary functions at random points: results missing the value & ---
& 0 of 1\,600\,000 & --- & --- & --- & 34\,\% (20\,\% with
\code{rounded\_}\allowbreak\code{transc\_std}) \\
\hline
\code{x + y} & 137 ns & 2.2 ns & 5.1 ns & 18 ns & 14 ns & 12 ns \\
\code{x * y} & 171 ns & 2.3 ns & 12 ns & 21 ns & 15 ns & 21 ns \\
\code{sin(x)} & 5.5 $\mu$s & 43 ns & 33 ns & 41 ns & 115 ns & 71 ns \\
\code{log(x)} & 1.9 $\mu$s & 21 ns & 27 ns & 39 ns & 10 ns & 30 ns \\
\code{pown(x, 3)} & 240 ns & 5.1 ns & 18 ns & 145 ns & 35 ns & 20 ns \\
Shekel 5 & 9.0 $\mu$s & 108 ns & 397 ns & 1.7 $\mu$s & 1.1 $\mu$s
& 920 ns \\
\hline
\end{tabular}

\end{table}

\begin{itemize}
\item[$\bullet$] \textbf{libieeep1788} gives the result of IEEE 1788 in
every case it can compute, as tightly as possible, and is 17 to 135 times
slower than GAOL.
\item[$\bullet$] \textbf{GAOL} is the fastest on the arithmetic, the
integer power and most formulas. Under the names of the standard,
\code{gaol\_ieee1788}, it gives IEEE 1788's result in every special case;
under its own, \code{gaol}, in all but 8, all from its hybrid \code{pow},
which takes \code{pown} for integer exponents, a negative base included. Its
elementary functions, CORE-MATH's, give the tightest bounds; its \code{exp}
and \code{log} are faster than filib++'s, its \code{sin} and \code{cos}
slower.
\item[$\bullet$] \textbf{filib++} is the fastest on \code{sin} and
\code{cos}, 1.3 to 1.6 times as fast as GAOL, but 5.3 times as slow on
$\times$, 3.5 times on $\div$ and 2.3 times on $+$ and $-$, and its bounds
of elementary functions and real powers are up to 36 doubles wider than the
tightest. Its extended mode gives other results than IEEE 1788 wherever an
infinity or a division by zero is involved, much as Solaris Studio:
\code{interval(+}$\infty$\code{)} is $\itv{\mathrm{MAX}}{+\infty}$ and
$\itv{0}{0} \times \itv{1}{+\infty}$ is $\itv{-\infty}{+\infty}$; its
\code{operator>>} rounds the bounds it reads to nearest, and its
comparisons with the empty set and the infinities differ from IEEE 1788's.
\item[$\bullet$] \textbf{PROFIL/BIAS} is a library of the interval
arithmetic before IEEE 1788: it has no empty set, and an operation whose
argument leaves the domain of the function, a division by an interval
containing zero included, prints an error and aborts the program. Its
exponential, logarithm and real power, taken from the libm, are the fastest
of the six libraries (9.9, 10.2 and 34 ns), and so is its square root (3.7 ns); its
arithmetic is 6.4 times slower than GAOL's on $+$ and $-$, its \code{sin}
and \code{cos} are the slowest after libieeep1788's (115 and 131 ns), and
its integer power, computed as $\exp(n \log x)$, is $1.1\cdot10^{-5}$ wider
than the tightest on average.
\item[$\bullet$] \textbf{Boost.Interval}, the library of Boost, is the
interval library C++ programmers find first. Its arithmetic gave the
tightest intervals in the benchmark, but its elementary functions are the C
library's, called with the rounding direction set downward or upward, which
glibc does not always honour: at the points of the test of enclosure,
34\,\% of its results miss the value of the function with glibc 2.43, 20\,\%
with its other policy, and \code{atan(}$\itv{-\infty}{+\infty}$\code{)}
misses $\pi/2$. With its default checking it has no empty set: an operation
that would create one throws an exception. It is 5.6 times slower than GAOL
on $+$ and $-$, 1.8 to 11 times on $\times$, $\div$ and the formulas, 3.5
times on the square root, 3.9 times on the integer power and 1.2 to 1.6
times on the other elementary functions.
\item[$\bullet$] \textbf{Solaris Studio} is 1.0 to 1.3 times as fast as GAOL
on \code{sin} and \code{cos}, but slow on integer powers and squares. Its
containment sets give other results than IEEE 1788 wherever an infinity, a
division by zero or an invalid argument is involved:
\code{interval(+}$\infty$\code{)} is $\itv{\mathrm{MAX}}{+\infty}$,
$\itv{1}{2} / \itv{0}{1}$ is $\itv{-\infty}{+\infty}$ and
\code{interval(2, 1)} is $\itv{-\infty}{+\infty}$.
\end{itemize}

\subsection{Performance}
\label{sec:computing-performance}
\cindex{benchmark of the libraries}%

One million operations of each kind are computed by each library on
ordinary intervals: bounded, of widths from $10^{-6}$ to 1, drawn once and
read from the same file by every program. Each operation is timed on the
whole array, its results being stored in another one; the time per
operation is the best time divided by the million. The same operations on
doubles, computed at the midpoints of the intervals without any rounding,
give the cost of the arithmetic itself. Table~\ref{tab:computing-bench-ops}
lists the operations: $a$ and $b$ are intervals centred in $\itv{-10}{10}$,
$p$ in $\itv{1}{10}$, $e$ in $\itv{0.5}{2.5}$, the arguments of Shekel 5 in
$\itv{0}{10}$; their widths are $10^u$, $u$ uniform in $\itv{-6}{0}$
($\itv{-6}{-1}$ for $e$).

\begin{table}[htbp]
\centering\footnotesize
\caption{The operations of the benchmark of \file{doc/compare/performance.md}.} \label{tab:computing-bench-ops}
\begin{tabular}{|l|L{10.5cm}|}
\hline
\textbf{Operation} & \textbf{Computes} \\
\hline
\code{add} & $a + b$ \\
\code{sub} & $a - b$ \\
\code{mul} & $a \times b$ \\
\code{div} & $a / p$ \\
\code{sqr} & $a^2$ (\code{sqr}, \code{square} in Boost.Interval,
\code{x**2} in Fortran) \\
\code{sqrt} & $\sqrt{p}$ \\
\code{exp} & $\exp(a)$ \\
\code{log} & $\log(p)$ \\
\code{sin} & $\sin(a)$ \\
\code{cos} & $\cos(a)$ \\
\code{pow\_int} & $a^3$ (\code{pow(x, int)} in GAOL and Boost.Interval,
\code{pown}, \code{power(x, int)} in filib++, \code{x**3}) \\
\code{pow\_real} & $p^e$ (\code{pow(x, y)}, \code{x**y}; $\exp(e \log p)$
with Boost.Interval, which has no real power) \\
\code{line\_arith} & $(a + b)(a - b) / p$ \\
\code{line\_trig} & $\sin(a) \cos(b) + a^2$ \\
\code{line\_pow} & $\sqrt{p} \cdot a^3 - \exp(b / p)$ \\
\code{shekel5} & Shekel 5: $-\sum_i 1 / (\sum_j (x_j - a_{ij})^2 + c_i)$,
5 terms, 4 variables \\
\code{block5} & five lines: $t_1 = ab + p$; $t_2 = \sin(t_1) \cos(b)$;
$t_3 = a^2 + t_2/p$; $t_4 = \exp(t_2) - b^3$; $t_3 t_4 + \sqrt{p}$ \\
\hline
\end{tabular}

\end{table}

Everything was built and run again with Clang 21.1, \option{-O3 -mfma} and
the flags of interval arithmetic (\pxref{sec:compiling}), so that one
compiler answers for every C and C++ library: GAOL v5 (V5.0.0) built by
CMake in Release with \option{GAOL\_FMA}, the sources of CORE-MATH compiled
into it and the rounding direction not preserved, its default, and called
under the names of IEEE 1788-2015, \code{gaol\_ieee1788}, as the other
libraries of the standard; the same GAOL built with
\option{GAOL\_PRESERVE\_ROUNDING}, which takes the AVX-512 path
(Section~\ref{sec:computing-avx512-rule}), in a column of the table of the
times only; GAOL 4.2.3, the last version of Fr\'ed\'eric Goualard, built
by its \cmd{configure} from his repository (cd0ee1a), static, the rounding
direction not preserved, its verbose mode off, with mathlib 2.1.1 built by
Clang 21.1 with \option{-O3 -mfma -ffp-contract=off}, and called under its
own names; libieeep1788 with MPFR 4.2.1 and GMP 6.3.0; filib++ 3.0.2.2, the
archive IBEX distributes, as
\code{interval<double, native\_switched, i\_mode\_extended\_flag>};
PROFIL/BIAS 2.0.8 in its \code{x86-64-Linux-compat-gcc} configuration;
Boost.Interval of Boost 1.92.0 with its default policies and
\code{rounded\_transc\_opp} for the elementary functions, glibc 2.43. Only
Solaris Studio keeps its own compiler, Sun Fortran 95 8.7 with
\option{-O3 -xia}, by nature. Every library was measured in the same run, on 17 October 2026, and GAOL v5 with the rounding direction preserved again on 18 October, on a laptop with an Intel i7-1185G7 at 3.00 GHz (pinned by
\cmd{taskset -c 2}), Linux 7.0.0-38-generic, Ubuntu 26.04.1; each time is
the best of the runs of three rounds, each program being run in turn with
the others. libieeep1788's times, from a single run of each operation in
each round, are the least steady. Solaris Studio's and PROFIL/BIAS's
operations are calls into libraries, as GAOL v5's \code{+}, \code{-},
\code{*}, \code{/} and \code{sqr} are, the bounds in registers, where the
compiler inlines GAOL 4.2.3's, filib++'s and Boost.Interval's.

\paragraph{What GAOL v5 changes.}
GAOL v5 is compared with GAOL 4.2.3, which it continues and which bounds
the elementary functions with mathlib (\pxref{sec:math-libraries}).
Table~\ref{tab:computing-v5-v4} gives the operations where they differ
most. The elementary functions are faster:\newinvfive{} two things go at
once with mathlib, the outward move of each bound, and the two changes of
rounding direction each function made, to nearest before mathlib and back
upward after. The integer power is computed from exact products, both
bounds at once in an SSE2 register, with an error bound proved beforehand.
The bounds are tighter, not looser: GAOL v5's results are no
wider than libieeep1788's, which computes every bound with MPFR, for every
one of the seventeen operations (excess $0.0$), where GAOL 4.2.3 is wider
on the square root, \code{exp}, \code{log}, \code{sin}, \code{cos}, both
powers and the lines that use them.

\begin{table}[htbp]
\centering\footnotesize
\caption[GAOL v5 against GAOL 4.2.3, in nanoseconds per operation]{GAOL v5 against GAOL 4.2.3, in nanoseconds per operation (Intel i7-1185G7, Clang 21.1, 17 October 2026).} \label{tab:computing-v5-v4}
\begin{tabular}{|l|r|r|l|}
\hline
\textbf{Operation} & \textbf{GAOL v5} & \textbf{GAOL 4.2.3} & \textbf{Ratio} \\
\hline
\multicolumn{4}{|l|}{\emph{The elementary functions and the powers}} \\
\hline
\code{log} & 20.9 ns & 46.8 ns & 2.2 times faster \\
\code{exp} & 18.3 & 35.1 & 1.9 \\
\code{pow(x, y)} & 48.2 & 89.9 & 1.9 \\
\code{pow(x, 3)} & 5.12 & 9.07 & 1.8 \\
line of powers & 38.1 & 61.0 & 1.6 \\
\code{sin} & 43.2 & 68.0 & 1.6 \\
five-line block & 155 & 229 & 1.5 \\
line of sin and cos & 102 & 148 & 1.5 \\
\code{cos} & 54.0 & 62.8 & 1.2 \\
\hline
\multicolumn{4}{|l|}{\emph{The arithmetic}} \\
\hline
\code{mul} & 2.30 ns & 10.0 ns & 4.4 times faster \\
\code{line\_arith} & 5.61 & 18.7 & 3.3 \\
\code{sqr} & 1.87 & 5.81 & 3.1 \\
\code{div} & 2.52 & 7.50 & 3.0 \\
\code{shekel5} & 108 & 189 & 1.7 \\
\hline
\end{tabular}

\end{table}

The arithmetic is faster than GAOL 4.2.3's, which is what the operators of
GAOL v5 are written for (Section~\ref{sec:kernels}). The sums and the
differences, 2.2 ns against 2.4, and the square root, 6.3 against 6.3, take
what they took: a sum is one instruction, and what it waits for is memory,
the benchmark reading two arrays of 16 MB and writing a third, and the
square root is the tightest interval whatever the rounding of the C
library's \code{sqrt}, where GAOL 4.2.3's is wider (excess
$1.4\cdot10^{-14}$). The check of the direction is what a program needs that
changes the rounding direction: the benchmark sums the results rounding to
nearest, and GAOL 4.2.3's own \code{midpoint()} leaves the direction to
nearest. Left so, the operations of GAOL 4.2.3 timed next gave bounds
narrower than libieeep1788's, which did not enclose the exact results; the
benchmark now sets the direction back after its sums. The bounds are the
tightest, as before: the zero bound of a product or of a quotient may have
the other sign (\pxref{sec:zero-sign}).

\paragraph{Preserving the rounding direction.}
The column of Table~\ref{tab:computing-times} where the direction is
preserved (\option{GAOL\_PRESERVE\_ROUNDING} on, with the AVX-512 path)
leaves the sums and the differences as they are, 2.2 ns, and costs 1.5
times on the products and 1.4 times on the quotients (3.4 and 3.6 ns), 1.1 times on
\code{sqr} (2.0 ns against 1.9) and 1.2 times on \code{pow(x, 3)} (6.2 ns
against 5.1): these operations take the AVX-512 instructions, whose rounding
direction is in the instruction, and set none, so that nothing is saved and
restored. The square root, 3.75 ns, is faster than without the preservation
(6.3 ns): the instruction rounds each root in its direction, where the SSE2
path checks the root of the C library with a product. The elementary
functions read the control register of the direction, set it upward and
write it back, once each, and take 1.0 to 1.3 times longer (\code{exp}
23.2 ns against 18.3, \code{log} 20.4 against 20.9, \code{sin}, \code{cos}
and \code{pow(x, y)} 1.1 times). They paid more, as \code{sqr} and the
integer powers did, before the kernels of Section~\ref{sec:kernels}: the two
writes of the control register, each one waiting for the operations in
flight, took 6 ns a \code{sqr} (8.3 ns, 4.4 times); with the C library's
\code{fegetround()} and \code{fesetround()} and the writes of the control
word of the x87 unit as well, \code{pow(x, 3)} took 14 ns more (19.0 ns, 3.7
times; 8.4 ns with the writes of the SSE register alone, where the machine
has no AVX-512) and the elementary functions 1.2 to 1.7 times longer.
Shekel 5 takes 147 ns against 108, and every operation of the table is
faster than with GAOL 4.2.3, which does not preserve the direction either
(Shekel 5: 189 ns). The results are the same: the sums of the midpoints and
of the widths of the million results are those of GAOL v5, to
$2\cdot10^{-16}$ relatively.

\begin{table}[htbp]
\centering\footnotesize
\setlength{\tabcolsep}{3pt}
\caption[Time per operation of the six libraries, in nanoseconds]{Time per operation, in nanoseconds: the same operations on doubles (double), libieeep1788 (p1788), GAOL v5 (v5), GAOL v5 with the rounding direction preserved and the AVX-512 path (v5 pres.), GAOL 4.2.3 (v4), filib++, Solaris Studio (Sun), PROFIL/BIAS (BIAS) and Boost.Interval (Boost). One million operations of each kind on the same intervals, the best of 3 rounds; Intel i7-1185G7, Clang 21.1.8, Ubuntu 26.04.1, 17 October 2026 (the preserved column on 18 October).} \label{tab:computing-times}
\begin{tabular}{|l|r|r|r|r|r|r|r|r|r|}
\hline
\textbf{Operation} & double & p1788 & v5 & v5 pres. & v4 & filib++
& Sun & BIAS & Boost \\
\hline
\code{add} & 0.96 & 137 & 2.23 & 2.19 & 2.38 & 5.10 & 17.6 & 14.2 & 12.4 \\
\code{sub} & 0.92 & 138 & 2.21 & 2.16 & 2.40 & 5.12 & 17.6 & 14.3 & 12.3 \\
\code{mul} & 0.90 & 171 & 2.30 & 3.42 & 10.0 & 12.1 & 21.3 & 14.7 & 21.5 \\
\code{div} & 0.91 & 171 & 2.52 & 3.60 & 7.50 & 8.70 & 20.3 & 14.4 & 18.9 \\
\code{sqr} & 0.59 & 88.4 & 1.87 & 2.04 & 5.81 & 8.25 & 55.5 & 18.5
& 16.2 \\
\code{sqrt} & 1.30 & 107 & 6.30 & 3.75 & 6.26 & 13.4 & 26.6 & 3.73
& 22.3 \\
\code{exp} & 3.44 & 1\,491 & 18.3 & 23.2 & 35.1 & 30.1 & 38.0 & 9.87
& 29.8 \\
\code{log} & 3.32 & 1\,919 & 20.9 & 20.4 & 46.8 & 26.6 & 38.5 & 10.2
& 29.9 \\
\code{sin} & 13.4 & 5\,472 & 43.2 & 46.8 & 68.0 & 33.4 & 41.5 & 115
& 71.2 \\
\code{cos} & 12.0 & 4\,890 & 54.0 & 57.6 & 62.8 & 33.9 & 42.7 & 131
& 65.7 \\
\code{pow\_int} & 13.0 & 240 & 5.12 & 6.18 & 9.07 & 17.8 & 145 & 35.1
& 19.9 \\
\code{pow\_real} & 10.3 & 6\,517 & 48.2 & 51.4 & 89.9 & 64.5 & 125 & 33.6
& 73.2 \\
\code{line\_arith} & 1.14 & 655 & 5.61 & 8.24 & 18.7 & 27.4 & 61.6 & 63.1
& 61.7 \\
\code{line\_trig} & 25.8 & 10\,920 & 102 & 113 & 148 & 90.8 & 166 & 311
& 187 \\
\code{line\_pow} & 19.4 & 2\,401 & 38.1 & 43.7 & 61.0 & 81.9 & 253 & 92.6
& 115 \\
\code{shekel5} & 2.43 & 8\,993 & 108 & 147 & 189 & 397 & 1\,720 & 1\,093
& 920 \\
\code{block5} & 44.5 & 13\,823 & 155 & 171 & 229 & 197 & 464 & 427 & 343 \\
\hline
\end{tabular}

\end{table}

\paragraph{What the timings show.}
On the arithmetic, GAOL is the fastest on the four operations: $+$ and $-$
take GAOL v5 2.2 ns (2.4 for GAOL 4.2.3), 2.3 times as fast as filib++
(5.1 ns) and 6.4 to 8.0 times as fast as PROFIL/BIAS (14 ns) and Solaris
Studio (17.6 ns); $\times$ and $\div$ take GAOL v5 2.3 and 2.5 ns (10.0 and
7.5 for GAOL 4.2.3), against 12.1 and 8.7 ns for filib++, 14.7 and 14.4 ns
for PROFIL/BIAS, 21 and 20 ns for Solaris Studio and 21 and 19 ns for
Boost.Interval. An addition of intervals costs GAOL v5 2.3 times an
addition of doubles: one vector addition for both bounds, and the check that the rounding
direction is upward, where filib++ sets the rounding direction and restores
it, and PROFIL/BIAS calls a function of BIAS that sets it downward, then
upward, then back to nearest, which is why its four arithmetic operations
all take about the same time. The integer power \code{pown(x, 3)} takes
GAOL v5 5.1 ns and GAOL 4.2.3 9.1 ns: filib++'s \code{power(x, 3)} takes
18 ns, Boost.Interval's 20 ns, PROFIL/BIAS's 35 ns and Solaris Studio's
\code{x**3} 145 ns; \code{sqr} takes GAOL v5 1.9 ns, against 8.2 for
filib++, 16 for Boost.Interval, 18 for PROFIL/BIAS and 56 for Solaris
Studio.

\medskip
On the elementary functions, PROFIL/BIAS is the fastest on the square root
(3.7 ns), \code{exp} (9.9 ns), \code{log} (10.2 ns) and the real power
\code{pow(x, y)} (34 ns), which it computes from the libm of the system,
moved outward, without correct rounding: its bounds are among the widest,
and its \code{sin} and \code{cos}, computed by its own argument reduction,
are the slowest after libieeep1788's, 115 and 131 ns. GAOL v5 comes next on
\code{exp} and \code{log}, 18 and 21 ns, faster than Boost.Interval (30 ns
each), filib++ (30 and 27 ns) and Solaris Studio (38 and 39 ns), with
correctly rounded values, where the others are not. filib++ is the fastest
on \code{sin} and \code{cos} (33.4 and 33.9 ns) with its own polynomials,
then Solaris Studio (41 and 43 ns), then GAOL v5 (43 and 54 ns): GAOL pays
there for correct rounding over the whole range, and for dividing the bounds
by an interval enclosing $\pi$ to find where the function is monotonic,
asking CORE-MATH for the signs of the derivative where the division cannot
tell. GAOL v5's \code{sin} became faster than its \code{cos} when CORE-MATH
rewrote \file{sin.c}. On \code{pow(x, y)} GAOL v5 takes 48 ns, against 65
for filib++, 73 for Boost.Interval, 90 for GAOL 4.2.3 and 125 for Solaris
Studio.

\medskip
On the formulas, GAOL is the fastest on the arithmetic line (5.6 ns for GAOL
v5 and 18.7 for GAOL 4.2.3, against 27 for filib++ and 62 to 63 for Solaris
Studio, PROFIL/BIAS and Boost.Interval), on the line of powers (38 and
61 ns, against 82 for filib++, 93 for PROFIL/BIAS and 253 for Solaris
Studio) and on Shekel 5 (108 and 189 ns, against 397 for filib++, 920 for
Boost.Interval, 1\,093 for PROFIL/BIAS and 1\,720 for Solaris Studio: Shekel
5 is 20 squares, 45 additions and subtractions and 5 divisions, and its
squares cost Solaris Studio 56 ns each and PROFIL/BIAS 18). GAOL v5 alone is
the fastest on the five-line block (155 ns, against 197 for filib++, 229 for
GAOL 4.2.3, 427 for PROFIL/BIAS and 464 for Solaris Studio). filib++ is the
fastest on the line of sin and cos (90.8 ns, against 102 for GAOL v5, 148
for GAOL 4.2.3, 166 for Solaris Studio and 311 for PROFIL/BIAS), where its
\code{sin} and \code{cos} weigh most.

\medskip
libieeep1788 is 17 to 135 times slower than GAOL v5: every bound is an MPFR
computation, about 140 ns for an addition and 5 to 7 $\mu$s for \code{sin},
\code{cos} and the real power. Its own documentation warns that its focus is
correctness, not speed. GAOL v5 gives the same bounds as it does, tightest
everywhere, 17 to 135 times faster. Boost.Interval is slower than GAOL v5 on
every operation: 5.6 times on $+$ and $-$, 3.5 to 9.3 times on $\times$,
$\div$, the square, the square root and the integer power, 1.2 to 1.6 times
on \code{exp}, \code{log}, \code{sin}, \code{cos} and the real power, 11
times on the arithmetic line and 8.5 times on Shekel 5. Each of its
operations reads the rounding direction (\code{fegetround}), sets it upward
(\code{fesetround}) and, at its end, sets the one it read back: two changes
of direction where GAOL checks the direction and changes nothing. Its
elementary functions add a change to the downward direction and back for the
lower bound, and its \code{sin} goes through the reduction of
$\cos(x - \pi/2)$ by an interval enclosing $2\pi$. Its bounds are those of
the C library, which are not correctly rounded, and of which 20 to 34\,\%
miss the value of the function with glibc 2.43
(Section~\ref{sec:computing-enclosure}).

\paragraph{The same results.}
The sums of the midpoints of the million results agree to $4\cdot10^{-15}$
relatively, apart from PROFIL/BIAS's and the sines of GAOL 4.2.3 and
Boost.Interval, $1.4\cdot10^{-13}$ and $1.3\cdot10^{-13}$ apart, their sum
being small, 373 for a million values, which makes $5\cdot10^{-17}$ per
result. The arithmetic operations give the tightest intervals in every
library, and so do the square roots, except those of GAOL 4.2.3, filib++ and
PROFIL/BIAS. Of the elementary functions, only GAOL v5's are the tightest on
every operation: on average the others are wider than libieeep1788's by up
to $5.9\cdot10^{-14}$ relatively with GAOL 4.2.3 (the real power, computed
as $\exp(y \log x)$), $2.1\cdot10^{-14}$ to $2.4\cdot10^{-13}$ with filib++ ($1.4\cdot10^{-13}$ for \code{log}, $2.4\cdot10^{-13}$ for the real power), $9.0\cdot10^{-15}$ to $1.2\cdot10^{-13}$ with PROFIL/BIAS, and at most $4.1\cdot10^{-15}$ with Solaris Studio; with Boost.Interval, by
$1.4\cdot10^{-14}$ and $2.5\cdot10^{-14}$ on \code{cos} and \code{sin} and
$3.7\cdot10^{-14}$ on the real power, its \code{exp} and \code{log} being as
tight as libieeep1788's on average. PROFIL/BIAS's integer power is far
wider, $1.1\cdot10^{-5}$ on average: it computes \code{Power(x, n)} as
$\exp(n \log x)$, where the others multiply.

\begin{Rem}
The arithmetic results of PROFIL/BIAS have midpoints that differ by
$8\cdot10^{-15}$ to $4\cdot10^{-14}$ relatively while their widths are the
tightest: the intervals are the same, but its \code{Mid} computes
$\inf + (\sup - \inf)/2$ rounded upward, one double above the midpoint for
4\,\% of the sums of the benchmark, where the others round to nearest. In
the same way, the sums of the widths of Boost.Interval's results of
\code{exp} and \code{log} equal libieeep1788's, while its upper bound of
\code{exp} is one double too low in 2 to 3\,\% of the results of the test of
enclosure: a sum of a million widths cannot show it, and only a check of
each bound against the exact value finds it.
\end{Rem}

\paragraph{Long lines of arithmetic operators.}
The operations above are timed one at a time. A program that computes with
intervals writes lines of them, an operator after another, where what an
operation hands to the next matters as much as what it computes. The lines
of Table~\ref{tab:computing-lines} are expressions of $+$, $-$, $\times$
and $\div$ alone, of 7 to 19 operations between 8 intervals of the same
million, evaluated for each of them, by every library but Solaris Studio,
which has no program for them. The seven first take intervals of both
signs, at random; the last three, \code{pos\_sum16}, \code{pos\_mul8} and
\code{pos\_quot8}, only $p$ and $e$, which are positive, so that the
branches of an implementation that tells the signs apart, as GAOL 4.2.3's
does, are taken the same way from one evaluation to the next. The chains
give the next operation the result of the previous one, the trees
(\code{mul8\_tree}, \code{mix\_tree}) several independent ones.
Table~\ref{tab:computing-lines-x86} gives the times on the x86 machine of
the benchmark, Table~\ref{tab:computing-lines-m2} on an Apple M2.

\begin{table}[htbp]
\centering\footnotesize
\caption[The long lines of arithmetic operators between 8 intervals]{The long lines of arithmetic operators between 8 intervals ($a$, $b$, $p$, $e$ as in Table~\ref{tab:computing-bench-ops}, $s_1$ to $s_4$ in $\itv{0}{10}$).} \label{tab:computing-lines}
\begin{tabular}{|l|L{9.0cm}|r|}
\hline
\textbf{Line} & \textbf{Computes} & \textbf{Operations} \\
\hline
\code{sum16} & $a + b + p + e + s_1 + s_2 + s_3 + s_4 + a + b + p + e + s_1
+ s_2 + s_3 + s_4$ & 15 \\
\code{diff16} & $a - b + p - e + s_1 - s_2 + s_3 - s_4 - a + b - p + e - s_1
+ s_2 - s_3 + s_4$ & 15 \\
\code{mul8} & $a \cdot b \cdot p \cdot e \cdot s_1 \cdot s_2 \cdot s_3 \cdot
s_4$ & 7 \\
\code{mul8\_tree} & $((a \cdot b)(p \cdot e))((s_1 \cdot s_2)(s_3 \cdot
s_4))$ & 7 \\
\code{quot8} & $a / p / e / p / e / p / e / p$ & 7 \\
\code{mix13} & $((((((a + b) \cdot p - e) \cdot s_1 + s_2) \cdot p - s_3) /
e + a) \cdot b - s_4) / p + e$ & 13 \\
\code{mix\_tree} & $((a + b)(p - e) + (s_1 - s_2)(s_3 + s_4)) / (p + e) +
((a - b)(s_1 + s_2) - (p \cdot e) / (p + e)) / (e + p)$ & 19 \\
\code{pos\_sum16} & $p + e + p + e + p + e + p + e + p + e + p + e + p + e
+ p + e$ & 15 \\
\code{pos\_mul8} & $p \cdot e \cdot p \cdot e \cdot p \cdot e \cdot p \cdot
e$ & 7 \\
\code{pos\_quot8} & $p / e / p / e / p / e / p / e$ & 7 \\
\hline
\end{tabular}

\end{table}

\begin{table}[htbp]
\centering\footnotesize
\setlength{\tabcolsep}{3pt}
\caption[Time per evaluation of a line, in nanoseconds]{Time per evaluation of a line, in nanoseconds, with the abbreviations of Table~\ref{tab:computing-times} (one million evaluations of each line, the best of 3 rounds; Intel i7-1185G7, Clang 21.1.8, 17 October 2026).} \label{tab:computing-lines-x86}
\begin{tabular}{|l|r|r|r|r|r|r|r|}
\hline
\textbf{Line} & double & p1788 & v5 & v4 & filib++ & BIAS & Boost \\
\hline
\code{sum16} & 2.45 & 2\,001 & 22.8 & 27.4 & 76.5 & 216 & 179 \\
\code{diff16} & 2.44 & 2\,091 & 23.3 & 27.1 & 76.5 & 215 & 177 \\
\code{mul8} & 2.42 & 1\,166 & 17.9 & 30.9 & 48.6 & 104 & 111 \\
\code{mul8\_tree} & 2.37 & 1\,168 & 15.0 & 29.4 & 50.5 & 105 & 111 \\
\code{quot8} & 3.36 & 1\,178 & 20.6 & 34.0 & 48.6 & 102 & 112 \\
\code{mix13} & 2.38 & 2\,068 & 23.0 & 47.6 & 80.7 & 192 & 185 \\
\code{mix\_tree} & 2.38 & 2\,985 & 34.7 & 64.7 & 116 & 278 & 264 \\
\code{pos\_sum16} & 1.03 & 1\,930 & 23.1 & 24.5 & 73.5 & 214 & 179 \\
\code{pos\_mul8} & 0.89 & 1\,139 & 17.8 & 18.9 & 40.2 & 100 & 105 \\
\code{pos\_quot8} & 3.27 & 1\,160 & 20.2 & 37.6 & 44.3 & 101 & 109 \\
\hline
\end{tabular}

\end{table}

GAOL v5 is the fastest of the libraries on every line, and faster than GAOL
4.2.3 on all ten, 1.06 to 2.1 times as fast: a line of 15 additions takes
22.8 ns, 1.5 ns an addition, against 27.4 ns (1.8 ns), and a line of 7
products 17.9 ns (2.6 ns a product) against 30.9 ns (4.4 ns); the quotients, 20.6 ns for 7 (2.9 ns each) against 34.0 ns (4.9 ns), 1.7 times as fast, and
the lines that mix the four operations, \code{mix13} and \code{mix\_tree},
23.0 and 34.7 ns against 47.6 and 64.7 ns, 2.1 and 1.9 times. The three
lines of intervals of one sign are where GAOL 4.2.3 comes nearest, on the
additions and the products, 1.06 times as fast as it, 23.1 ns against
24.5 ns and 17.8 ns against 18.9 ns: with the signs the same from one
evaluation to the next, its branches are predicted, and what it still pays
is the bounds that go through memory between two operations; its quotients,
on the other hand, take 1.9 times as long (37.6 ns against 20.2 ns). An
operation of a line costs GAOL v5 less than the same operation timed alone,
1.5 ns an addition against 2.2: the two doubles of the result stay in a
register for the next operation, and the benchmark of a single operation,
which reads two arrays and writes a third for each, waits for the memory,
where a line reads eight arrays at most for its 7 to 19 operations. Against
the other libraries, filib++ takes 2.2 to 3.5 times as long as GAOL v5 on
the ten lines, Boost.Interval 5.4 to 8.0 times, PROFIL/BIAS 5.0 to 9.5 times
and libieeep1788 57 to 90 times. The bounds are the tightest: the mean
widths of the results of GAOL v5 are those of libieeep1788 to
$1.1\cdot10^{-16}$ relatively on the lines of additions and exactly on the
others, where GAOL 4.2.3's differ from them by $2.3\cdot10^{-16}$ on the
lines of quotients.

\begin{table}[htbp]
\centering\footnotesize
\caption[The same lines on a 64-bit ARM processor (Apple M2)]{The same lines on a 64-bit ARM processor, in nanoseconds per evaluation: GAOL v5 with the NEON intervals (v5), the same with the rounding direction preserved (v5 pres.) and GAOL 4.2.3 with the FPU intervals (v4). Apple M2, Apple clang 21.0.0, macOS 27.0.1, 6 November 2026, the best of 3 rounds, the processor not pinned.} \label{tab:computing-lines-m2}
\begin{tabular}{|l|r|r|r|r|r|r|}
\hline
\textbf{Line} & \textbf{Ops} & v5 & v5 pres. & v4 & v5 pres. / v5
& v4 / v5 \\
\hline
\code{sum16} & 15 & 9.11 & 385 & 41.9 & 42.3 & 4.6 \\
\code{diff16} & 15 & 9.30 & 390 & 77.6 & 42.0 & 8.3 \\
\code{mul8} & 7 & 10.5 & 187 & 28.3 & 17.8 & 2.7 \\
\code{mul8\_tree} & 7 & 9.95 & 190 & 25.0 & 19.1 & 2.5 \\
\code{quot8} & 7 & 17.6 & 229 & 37.3 & 13.0 & 2.1 \\
\code{mix13} & 13 & 14.3 & 353 & 61.6 & 24.7 & 4.3 \\
\code{mix\_tree} & 19 & 18.5 & 516 & 82.2 & 27.9 & 4.4 \\
\code{pos\_sum16} & 15 & 8.85 & 385 & 39.7 & 43.5 & 4.5 \\
\code{pos\_mul8} & 7 & 9.58 & 187 & 16.1 & 19.6 & 1.7 \\
\code{pos\_quot8} & 7 & 17.6 & 232 & 46.5 & 13.2 & 2.6 \\
\hline
\end{tabular}

\end{table}

On the Apple M2, GAOL v5 takes 8.85 to 18.5 ns for a line of 7 to 19
operations, 0.6 to 2.5 ns an operation, where GAOL 4.2.3, which has no instructions of 64-bit ARM to compute its FPU intervals with, takes 1.7 to
8.3 times as long (its \cmd{configure} was generated again by
\cmd{autoreconf} and told \option{--build=arm-apple-darwin}, and mathlib
2.1.1 configured with \option{--host=aarch64-unknown-linux-gnu}; no line
calls it). With the rounding direction preserved the same lines take 187 to
516 ns, 25.7 to 33.1 ns an operation, 13.0 to 43.5 times as long as without:
the program computes to nearest, so each operation sets the direction upward
and puts nearest back, two \code{msr} of FPCR of 11 ns each that change it,
with the \code{mrs} of 0.8 ns that reads it. The sums of the midpoints and
of the widths of the three agree to $1.1\cdot10^{-16}$ relatively. The
machine was on mains power, the low power mode off, with a load average of
1.7 at the start and 3.0 at the end (the run took 65 s, the build of the
three programs included), and no process above 10\,\% of a core but the interface of
the machine in the readings taken just before and just after the run; the
three rounds of a line differ by 1.4\,\% at the median and by 6.4\,\% at
most, and the best of the three is kept.

\subsection{The special cases}
\label{sec:computing-special-cases}
\cindex{special cases!comparison}%

The special cases of GAOL's tests (\file{tests/*.cpp}) are computed by the six
libraries, GAOL v5 twice, under its own names and under those of IEEE
1788-2015: zeros, infinities and NaN as bounds or operands, empty sets,
divisions by intervals containing zero, arguments at the edges of the
domains, the powers \code{pow} and \code{pown}, the operators \code{+=},
\code{-=}, \code{*=}, \code{/=} with special doubles, the reading of
intervals from text, the numeric functions of intervals, the comparisons
and \code{atan2}. Each result is compared with the result IEEE 1788-2015
defines for the set-based flavor, the tightest interval of doubles, computed
with mpmath, and classed as the result of IEEE 1788, as an interval
enclosing it, wider, as another result (a set not enclosing it, a set where
IEEE 1788 has the empty set, another boolean, or an exception), or as not
applicable when the library has no such operation; IEEE 1788 defines no
result in two cases (a bare number read as an interval, and an n-th root of order 0). Signed zeros are not told apart. Each case is written
once, in \file{doc/compare/code/cases.py}, and translated for each library; the errors
of PROFIL/BIAS, which abort the program, and the exceptions of
Boost.Interval are caught. The counts are those of
Table~\ref{tab:computing-inshort}.

\medskip
libieeep1788 gives the result of IEEE 1788 in each of the 281 cases where
IEEE 1788 defines one and it has the operation; it has no n-th root, no
operators with doubles and no compound assignments. GAOL under the names of
the standard gives the result of IEEE 1788 in each of the 288 cases where
IEEE 1788 defines one.\newinvfive{} Under its own names it gives it in 280
cases, an interval enclosing it in 2, and something else in 6, all from the
hybrid \code{pow(x, y)} of \code{gaol}, which takes the integer power for an
integer exponent, where IEEE 1788's \code{pow} only takes the part of $x$ in
$\itv{0}{+\infty}$ (\pxref{sec:pow-namespaces}): \code{pow([-2], 2.0)} is
$\itv{4}{4}$, \code{pow([0], [0])} is $\itv{1}{1}$,
\code{pow([-1], [2\textasciicircum{}31-1])} is $\itv{-1}{-1}$,
\code{pow([-2, -1], [1e10])}, an integer beyond the ints, is
$\itv{1}{+\infty}$, and \code{pow([-2, 0], [-1])} is $\itv{-\infty}{-0.5}$,
where IEEE 1788 has $\emptyset$; \code{pow([-2, 3], 3)} is $\itv{-8}{27}$,
wider than $\itv{0}{27}$. In GAOL~4, the integer powers beyond the ints were
$\itv{-\infty}{+\infty}$, which made three more cases wider.\newinvfive{} GAOL reads a bare number, \code{textToInterval("0.1")},
as the interval enclosing it, an extension of the literals IEEE 1788 allows
(\pxref{sec:reading-intervals}). Its \code{sin}, \code{cos}, \code{tan},
\code{acos}, \code{acosh}, \code{atan2}, integer powers and n-th roots are
the tightest in every one of these cases, the powers being computed from exact products,
$x^{-n}$ included, and the roots proved with integer powers; across the
half-line $y = 0$, $x < 0$, where the angle jumps from $\pi$ to $-\pi$,
\code{atan2} is $\itv{-\pi}{\pi}$, as in libieeep1788.

\medskip
The other libraries depart from the standard each in its own way. filib++,
in its extended mode, takes the infinities for values, much as Solaris
Studio does: the intervals built from an infinity are
$\itv{\mathrm{MAX}}{+\infty}$, a product of a zero and an infinite bound is
$\itv{-\infty}{+\infty}$, \code{log([0])} is $\itv{-\infty}{-\mathrm{MAX}}$
and \code{atanh([1])} is $\itv{\mathrm{MAX}}{+\infty}$ where IEEE 1788 has
$\emptyset$; its \code{operator>>} reads only the form \code{[l, u]} and
rounds the bounds to nearest, so that \code{[0.1, 0.3]} does not enclose
$\itv{0.1}{0.3}$; its elementary functions are wider than the tightest by
up to 25 doubles and its real powers by up to 36. PROFIL/BIAS, a library of
1993 to 2009, has no empty set: an argument outside the domain of a
function, or a division by an interval containing 0, prints an error and
calls \code{abort()}; the intervals built from an infinity, a NaN or bounds
in the wrong order are not empty, \code{interval(2, 1)} being $\itv{1}{2}$;
its \code{sin} and \code{cos} are wider than the tightest by 4 to 14
doubles, its real powers by 3 to 14. Solaris Studio follows the containment
sets of Sun's interval arithmetic, where $1/0$ is $\{-\infty, +\infty\}$ and
$+\infty$ is a point: \code{interval(2, 1)} and NaN bounds give
$\itv{-\infty}{+\infty}$ rather than $\emptyset$, a division by an interval
containing zero is $\itv{-\infty}{+\infty}$, a single number read as an
interval carries an uncertainty of one unit of its last digit (\code{0.1}
gives $\itv{0}{0.2}$), and its powers are often one double wider than the
tightest. Boost.Interval, with its default checking, throws
\code{std::runtime\_error} wherever the empty set would appear, the empty
set given as an argument included; an infinity is a point, a product of a
zero and an infinity has a NaN bound, its \code{median} overflows for
$\itv{\mathrm{MAX}/2}{\mathrm{MAX}}$, and \code{atan([-}$\infty$\code{, +}%
$\infty$\code{])} is $\itv{-1.5707963267948966}{1.5707963267948966}$, which
does not hold the values of $\arctan$ near $\pm\pi/2$: Boost.Interval asks
the C library for $\arctan(+\infty)$ rounded upward, and glibc rounds it to
nearest, below $\pi/2$, whatever the rounding direction.

\subsection{The enclosure of the elementary functions}
\label{sec:computing-enclosure}
\cindex{enclosure test}%

An interval library promises one thing: the interval it returns contains
the exact result. The special cases check the promise at the edges of the
domains; the test of enclosure checks it in the middle of them. Each of 14
elementary functions, the square root and $x^3$ is computed on the point
intervals $\itv{x}{x}$ of 100\,000 random arguments $x$, half of them
ordinary, half at the edges of the doubles or of the domain, and each result
is compared with $f(x)$, computed by mpmath with 2000 bits (8000 where
$f(x)$ is that close to a double, as $\tanh(709)$ is to 1). A result
$\itv{l}{u}$ misses $f(x)$ when $l > f(x)$ or $u < f(x)$. Boost.Interval
computes its arithmetic in the rounding direction each bound needs, but its
elementary functions are those of the C library, called with the rounding
direction set downward for the lower bound and upward for the upper one:
they enclose $f(x)$ only if the C library honours that direction, which its
documentation says is rarely the case. GAOL v5, whose elementary functions
are CORE-MATH's, correctly rounded in each direction, is the control: none
of its results may miss. Both are compiled by Clang 21 with the flags of
interval arithmetic, Boost 1.92.0 with glibc 2.43 (Ubuntu 26.04.1), under
its two policies: \code{rounded\_transc\_std}, which sets the direction
downward, calls the function for the lower bound, sets it upward and calls
it for the upper one, and \code{rounded\_transc\_opp}, the policy of the
benchmark, which keeps the direction upward and sets it downward only around
the call for the lower bound.

\medskip
GAOL v5 encloses $f(x)$ in every one of the 1\,600\,000 results, under both
namespaces, and gives the tightest interval of doubles around it in all of
them (1\,232 results of $x^3$, where the cube is below $2^{-968}$, came from
products rounded outward and were one or two doubles wider, until these
corners were computed apart, scaled; \pxref{sec:accuracy-powers}). Boost.Interval
misses $f(x)$ in 34\,\% of the results with \code{rounded\_transc\_opp} and
in 20\,\% with \code{rounded\_transc\_std}. Its square root and its $x^3$,
which it computes itself with the rounding of the processor, never miss,
nor do its \code{asinh} and \code{atanh}; every other elementary function
misses, \code{log} once, the others in 0.05\,\% of the results or more, and
\code{asin}, \code{atan}, \code{sinh}, \code{tanh} and \code{acosh} in
100\,\% with \code{rounded\_transc\_opp}. Three causes were found. The
\code{sin}, \code{cos}, \code{tan} and \code{atan} of glibc ignore the
rounding direction: called downward, to nearest and upward, they give the
same double for every argument, which misses $f(x)$ unless $f(x)$ is that
double; Boost.Interval's \code{sin}, \code{cos} and \code{tan} miss less
often, 0.3 to 28\,\%, because they reduce the argument by subtracting a
multiple of an interval enclosing $\pi$ or $2\pi$, which widens most results
by tens to hundreds of doubles. The \code{exp}, \code{cosh}, \code{acos} and
\code{log} of glibc are computed in the direction asked, but not always on
the right side of $f(x)$: called upward, \code{exp} gives a double below
$f(x)$ for 2 to 3\,\% of the arguments, \code{cosh} for up to 7\,\%,
\code{acos} misses in 0.09 to 5\,\%, \code{log} once in the 100\,000. The
\code{asin}, \code{sinh} and \code{tanh} of glibc round the wrong way for
many arguments, the value computed downward being above the one computed
upward, as if the function were computed for $|x|$ and its sign changed
afterwards: with \code{rounded\_transc\_std} the result is then inverted,
its lower bound above its upper one, and the default checking does not see
an inverted interval as empty, since it holds that no interval is.

\begin{Rem}
\code{rounded\_transc\_opp} has a bug of its own. For \code{sin},
\code{tan}, \code{asin}, \code{atan}, \code{sinh}, \code{tanh} and
\code{acosh} it computes the lower bound as the opposite of the opposite of
$f(x)$ computed upward, that is $f(x)$ computed upward again: the lower
bound is the upper one. The trick of the opposite, which works for the
arithmetic, would need $f_{\downarrow}(x) = -f_{\uparrow}(-x)$, and only for
the odd functions, which \code{acosh} is not. Its results of \code{asin},
\code{sinh}, \code{tanh} and \code{acosh} are therefore single doubles,
which miss $f(x)$ whenever it is not a double. The results depend on the C
library: glibc 2.31, under Ubuntu 20.04, missed in 41\,\% and 27\,\% of the
results, \code{asinh}, \code{atanh} and \code{acosh} included.
\end{Rem}

%%--------------------------------------------------------------------------
\section{ITF1788 and the benchmark of Tang et al.}
\label{sec:test-suites}
%%--------------------------------------------------------------------------
\cindex{ITF1788}\cindex{benchmark of Tang et al.}%

Two test suites written outside GAOL were run on GAOL v5 under the names of
the standard, \code{gaol\_ieee1788}, with the default build (SSE2 intervals
on x86, the machine of the tables being an Intel i7-1185G7) and the FPU build. The tests of GAOL
itself are the subject of Chapter~\ref{chap:tests}.

\paragraph{ITF1788.}
ITF1788, the Interval Test Framework for IEEE 1788 of Marco Nehmeier,
Maximilian Kiesner and Oliver Heimlich
(\url{https://github.com/oheim/ITF1788}), holds test cases of the
operations of IEEE 1788-2015: those of libieeep1788, of C-XSC, of FI\_LIB
and of MPFI, which Oliver Heimlich converted, and his own for \ieee{atan2},
\ieee{absRev}, \ieee{powRev1}, \ieee{powRev2}, the constructors and the
exceptions. They are written in a format of its own, one case per line of
its \file{.itl} files, such as
\code{add [1.0,2.0] [3.0,4.0] = [4.0,6.0];}. ITF1788 converts them into test
programs for several libraries; its plugin for GAOL, written for GAOL~4,
calls GAOL by GAOL's own names and leaves out most of the operations GAOL
v5 adds. The generator of \file{doc/compare/itf1788/} reads the
\file{.itl} files itself and writes a C++ program per file, each case being
one call of the operation by the name \code{gaol\_ieee1788} gives it.

\medskip
The suite holds 9542 cases in 19 files. 7415 pass: GAOL v5 gives the result
the case expects, the tightest one, the 1347 cases of \ieee{pow}, the 163
of \ieee{pown}, the 172 of \ieee{mulRevToPair} without decorations, the 48
of \ieee{overlap} without decorations, 802 of the 804 cases of
\ieee{powRev1} and \ieee{powRev2}, and 84 of the 94 cases of \ieee{pownRev}
with an exponent $p \leq 0$ included. No bug of GAOL v5 was found. The 22
failures are of three kinds: 12 results that are valid but not the
tightest, from \ieee{pownRev}, 2 at the largest double and 10 with an
exponent $p < 0$: 11 of them 1 double beyond the tightest bound and accurate in the sense of IEEE 1788-2015, and one 2 doubles beyond, valid but not accurate, the standard requiring of
this operation a valid result and recommending an accurate one; 8 cases where
ITF1788 expects a result that is not the tightest, from the reverse
functions \ieee{sinRev}, \ieee{cosRev} and \ieee{tanRev} of libieeep1788's
tests and \ieee{powRev2} of Oliver Heimlich's, where GAOL v5's result is
narrower than the expected one and encloses the exact result, which mpmath
computes; and 2 cases where ITF1788 expects an interval for a literal
\code{[l, u]} with $l$ above $u$, \code{[1.0000000000000002,1.0000000000000001]}
and the same bounds in hexadecimal, where GAOL v5 compares the numbers
exactly and gives the empty set, as the standard asks of an implementation
that decides $l \leq u$. 418 cases are not provided: eight trigonometric
functions of FI\_LIB and MPFI that are not in the standard (\code{cot},
\code{sec}, \code{csc}, their hyperbolic versions, \code{acot} and
\code{acoth}: 403) and the reduction operations (15). 1687 cases are not
applicable: they check the decorated intervals and NaI, the ill-formed interval of the
standard (11.3), which GAOL v5 does not have (see Chapter~\ref{chap:ieee1788}). The default and the FPU builds give the
same result in every case. The deliberate differences of GAOL v5 from the
standard do not show here: the hybrid \code{pow} is \code{gaol::pow}, while
\code{gaol\_ieee1788::pow} is the \code{pow} of the standard, and the
constructors of intervals with an infinite bound give the empty set where
the cases expect it.

\paragraph{The benchmark of Tang et al.}
The benchmark of interval libraries of Xuan Tang, Zachary Ferguson, Teseo
Schneider, Denis Zorin, Shoaib Kamil and Daniele Panozzo, \emph{A
Cross-Platform Benchmark for Interval Computation Libraries}
\cite{Tang:PPAM2022}, compares filib
(its C version), filib++, Boost.Interval and BIAS on 132 expressions: 28 of
the authors (the four operations, \code{sqrt}, \code{exp}, \code{sin} and
\code{cos}, ten arithmetic expressions of 6 to 10 variables and ten random
expressions with \code{sqrt}, \code{exp}, \code{sin} and \code{cos}) and 104
expressions of FPBench, on the inputs of their preconditions. Each
expression is evaluated on random inputs, each input variable being the
interval of one double; a result is correct when it contains the exact
value of the expression at the input, which the benchmark checks by printing
each result as a query with exact rationals. The paper finds every library
correct on the arithmetic, Boost.Interval wrong on \code{exp}, \code{sin}
and \code{cos}, BIAS wrong on one random expression, and filib++'s
\code{native\_switched} mode wrong on four FPBench expressions; GAOL is
cited but not run. A patch of \file{doc/compare/tang2021/} adds GAOL v5 to
it, and a checker replaces the Mathematica script of the paper.

\medskip
On 100\,000 inputs of each of the 132 expressions, 13\,200\,000 results in
each of the default and FPU builds of GAOL (CMake Release by Clang 21.1.8,
Intel i7-1185G7, Ubuntu 26.04, 14 October 2026), every interval of GAOL
encloses the exact value of the expression the program computes, the eight
single operations give the tightest interval every time (800\,000 of
800\,000), and GAOL's intervals are narrower than those of filib's C
version, the paper's reference, at all but 1\,488 inputs, where they are as
wide. The paper's own criterion fails at 9\,884 inputs, in six
FPBench expressions: each has a decimal constant that is not a double, which
the benchmark's query writes as the decimal while the program computes with
the double. These are a defect of the benchmark, not of GAOL; no bug of GAOL
was found. The two builds give the same tables, line for line. The speed
test of the benchmark, 1000 inputs of each expression each evaluated 10\,000
times, on GAOL and filib's C version built by Clang 21 in the SSE2 build, on
the same machine and day, with a load of about 1.5, is indicative only: GAOL
takes less time than filib's C version on 93 of the 132 expressions, 0.36
to 0.66 times its time on the four operations, 0.40 to 0.61 on the ten
arithmetic expressions, 0.55 for the median expression, and 107.9 s in all
against 131.7 s; it takes more on the elementary functions, 1.54 times on
\code{exp}, 1.77 on \code{sqrt}, 1.98 on \code{sin} and 3.98 on \code{cos},
whose inputs the speed test draws among $\pm(2^{31} - 1)$, where GAOL's
bounds are the tightest and filib's far wider on large arguments.

%%--------------------------------------------------------------------------
\section{The verdicts on new paths}
\label{sec:new-paths}
%%--------------------------------------------------------------------------
\cindex{paths of GAOL!verdicts}\cindex{RISC-V}\cindex{SVE}%

The sections above measured what GAOL has. This section draws the verdicts
on the tracks that were considered for new paths, lists the cheap gains
that need no new instruction set, states what a new path has to show
before it is added, and ends with the things to watch and what was not
measured. Table~\ref{tab:computing-verdicts} gives the verdicts.

\begingroup\footnotesize
\begin{longtable}{|L{2.3cm}|L{2.3cm}|L{5.0cm}|L{3.4cm}|}
\caption{The verdicts on the tracks considered for a path of GAOL (from
\file{doc/simd.md}).}\label{tab:computing-verdicts}\\
\hline
\textbf{Track} & \textbf{Verdict} & \textbf{Why} & \textbf{What would
change it} \\
\hline
\endfirsthead
\multicolumn{4}{l}{\textit{Table~\thetable, continued}}\\*
\hline
\textbf{Track} & \textbf{Verdict} & \textbf{Why} & \textbf{What would
change it} \\
\hline
\endhead
\hline
\endfoot
\endlastfoot
SSE4.1, AVX, AVX2 (a 256-bit path) & no path & A batch of two intervals per
register needs an interface of arrays, and nothing makes one interval
faster; \code{ROUNDPD} rounds to an integer, not as an operation & An
interface of boxes or arrays of intervals, and a program bound by the
computation and not by the memory \\
\hline
FMA3 & in use (\code{GAOL\_FMA}) & It pays in the elementary functions and
the integer powers, none in the basic operations & The exact
transformations that would replace the directed rounding are another study
(the rounding direction left to nearest) \\
\hline
AVX-512 path & keep; the division and the square root are scalar
instructions with the embedded rounding, two for the two bounds & On Intel
the 512-bit division takes 16 cycles of throughput against 4 for the
128-bit one (uops.info), and Zen 5 has full-width units; the square root of
the path takes two scalar embedded roots (3.75 ns against 6.30 with SSE2),
\code{x / y} two scalar divisions, and so do \code{x / d} of a double
(2.8 ns against 3.7 for $d > 0$, 3.0 to 3.3 for $d < 0$, on an Intel
i7-1185G7 with Clang 21), the divisors with a zero (5.5 ns) and the negative
integer powers (\code{pow(x, -3)} 9.1 ns), which the two scalar divisions, in place of the 512-bit one, left within the
noise & A processor where the two scalar divisions lose to the 512-bit one
(not seen on the i7-1185G7; no other processor was measured, Zen 5
included) \\
\hline
AVX-512VL, AVX10.2 at 128 and 256 bits & nothing to build & No embedded
rounding at these widths & A new revision of the Intel manual restoring it
(revision 3.0 of AVX10.2 had it) \\
\hline
NEON (arm64) & keep & Done: 1.3 to 2.5 times faster than the FPU intervals on an M2, 3.2 times on a line of 13 operations & Hardware with
FPCR.AH set (below) \\
\hline
SVE, SVE2, SME & no path & The same FPCR, and 128-bit vectors on most of the
processors & An interface of arrays, on a 256 or 512-bit implementation; a
static rounding in the instruction set \\
\hline
ARMv7 & no vector path & NEON has no double & none \\
\hline
POWER (VSX), s390x (vector facility) & measure first & Two doubles per
register and a write of the control register, as NEON, but no hardware here
and no figure for the write & A machine to run the tests and the benchmark
on \\
\hline
RISC-V, scalar & measure on a board, then write it & The only other rounding
direction by instruction, by inline assembly & A board, and a check that its
cores do not slow the static modes \\
\hline
RISC-V vector, LoongArch, MIPS, WebAssembly & no path & No static rounding
for the vectors, no user, or the rounding to nearest only & An extension
with a static rounding; for WebAssembly, a lifting of its limit \\
\hline
\end{longtable}
\endgroup

\subsection{Cheap gains that need no new instruction set}

The candidates below were seen in the sources and the profiles; the first
and the third are done, and none needs a new path.

\begin{itemize}
\item[$\bullet$] \textbf{\code{pow(x, 2)} through \code{sqr}} (done).
\code{gaol\_pown} took 5.2 ns where \code{sqr} takes 1.3 on the M2 (NEON
intervals; 5.2 and 1.8 with the FPU ones, whose \code{sqr} has no branch on
the signs), and the integer powers are about 40\,\% of Goldstein-Price (6
million calls). It takes \code{sqr} now, for $n = 2$ in every form of the
integer power:\newinvfive{} 1.6 ns with the NEON intervals and 1.9 with the
FPU ones, the call into the library accounting for the 0.3 ns of NEON above
\code{sqr} (0.1 with the FPU ones; the medians of 7 passes alternated with
the library before the change, \file{tests/performance.cpp}, an Apple M2,
AppleClang 21). The other powers cost 0.3 to 0.5 ns more with the FPU
intervals for $n = 3, 4, 5$ and 10 ($+5\,\%$), and less with the NEON ones.
On 14 million cases (3 million intervals, the special bounds, the empty set
and the universe, in the four directions of the calling code and with
flush-to-zero set) the bounds of \code{pow(x, 2)} are the same as before,
the sign of a zero lower bound aside, which is that of \code{sqr(x)} now.
\item[$\bullet$] \textbf{\code{sqrt} on SSE2} is scalar (6.3 ns on the
i7-1185G7, 3.75 with the AVX-512 path and the direction preserved, 2.0 on
the M2 with NEON); the NEON kernel ports to \code{\_mm\_sqrt\_pd}.
\code{sqrt\_rel} follows.
\item[$\bullet$] \textbf{The integer powers on NEON} (done, 1.1 to 1.3 times).
The kernel of \code{pown}, in \file{gaol/core/gaol\_interval\_neon.cpp},
computes on two lanes with \code{vfmaq\_f64}, as the SSE2 one of
\file{gaol\_interval\_sse.cpp} does with \code{\_\_FMA\_\_}, where it was
two scalar lanes with \code{std::fma}; the generic part is in
\file{gaol\_interval\_core.cpp}, and each kind of intervals defines the
operations of its lanes as a class.\newinvfive{} The bounds are the same bit
for bit (4 million powers hashed, and \file{tests/tools/pow/diffpow.py} on
70\,500 lines). The first version gained only 1.03 to 1.07 times, far below
the 1.5 times that was hoped for: \code{pow(x, 3)} took 5.7 ns where it took
6.1 (5 passes alternated by \file{.github/scripts/performance-ab.sh}, an
Apple M2, AppleClang 21). The kernel was not to blame: it is some 8 of the
125 instructions of a call, and its chain of dependent products is no
shorter than that of two scalar lanes (NEON has no vector \code{fnmsub}: each
fused multiply-subtract is a \code{vnegq} and a \code{vfmaq}). The rest was
the code before it, which Clang made longer: with \option{-frounding-math} it
turns the generic comparisons \code{vcgeq\_f64}, \code{vcleq\_f64} and
\code{vceqq\_f64} of the check of the range into scalar ones, 38
instructions where the same masks with \code{vcageq\_f64} and
\code{vcaleq\_f64}, on the absolute values, take 20, and it made the
permutation of the corners for a negative exponent (\code{vextq\_f64}) a
select through a general register. The check is now written with
\code{vcageq\_f64} and \code{vcaleq\_f64}, and the permutation keeps an empty
\code{asm} that leaves it as it is: a call of \code{pow(x, 3)} takes 102
instructions where it took 125. Without \option{GAOL\_PRESERVE\_ROUNDING}, it
takes 4.42 ns where it took 5.84 (1.32 times), and the other exponents from
$-7$ to 100 gain 1.11 to 1.33 times in independent calls (medians of 21
alternations of the old and the new code, an Apple M2, AppleClang 21). With
it, the 27 rows of the micro-benchmark take 0.95 to 0.99 of their time, a call
of \code{pow(x, 3)} taking 36.4 ns where it took 37.2: the preservation itself
adds some 31 ns to a call, whose cause was not looked for. Making the corners
from the stored register rather than from two stores and a load had not
changed the ratio of the first version.
\item[$\bullet$] \textbf{\code{mid()}} takes 8.5 to 12 ns on the M2 for an
addition at 0.9; a two-lane kernel would keep the rare branches.
\item[$\bullet$] \textbf{The argument reduction of \code{sin}, \code{cos}
and \code{tan}} (the quotients by $\pi$ and the \code{floor} and
\code{ceil}) is part of the code of GAOL around CORE-MATH, which is 42 to
66\,\% of the time of \code{sin} in the two programs that call it; no
profile says how much of that is the reduction.
\item[$\bullet$] \textbf{\code{div\_rel} and \code{\%=}} are trees of
branches; the NEON intervals have the cases of the function
\code{divide\_by\_zero\_containing()} as a table for the quotient.
\item[$\bullet$] \textbf{Not for SIMD.} \code{midpoint()} (29 ns on
the M2), \code{bisect()}, \code{chi()} and \code{hausdorff()} spend their
time in changes of the rounding direction, which is the question of
computing without changing it; the elementary functions are limited by
CORE-MATH, and a vector version would be another algorithm.
\end{itemize}

\subsection{What a new path has to show}

Before a path is added, with the rules of GAOL, it has to show:
\begin{enumerate}
\item its time against the path it replaces, on the processor and the
compiler that matter, with the spread, and the operation or the program
that needs it;
\item its bounds: the tightest ones, against references computed apart, in
the four rounding directions of the calling code, with and without
\option{GAOL\_PRESERVE\_ROUNDING}, the subnormals (the flush-to-zero modes
cleared), the NaN and the infinities, the empty set without an
invalid-operation exception, and the sign of a zero bound documented where
it differs;
\item the three build systems alike, a job of the continuous integration
that tells which path it took (\file{.github/scripts/check-intervals.sh},
\pxref{sec:ci}), and the processor it needs written with the platforms
(\pxref{sec:platforms}), with the path it falls back to.
\end{enumerate}

\subsection{Things to watch}

\begin{itemize}
\item[$\bullet$] \textbf{FPCR.AH.} With this bit of FEAT\_AFP (Armv8.7,
present on the Neoverse V3 of Graviton 5 and absent from the M2) set,
\code{FMAX} and \code{FMAXP} (the \code{vmaxq\_f64} and \code{vpmaxq\_f64}
of the NEON kernels) return the second operand for two zeros of different
signs, and when either is a NaN, in which case they also raise the
invalid-operation exception. GAOL clears FZ and FIZ and not AH. The kernels
never give a NaN to a maximum, and the zeros of the maxima of the product
differ only in sign. The square root used to depend on the order of the
operands of a maximum, of the interval with its left bound in its natural
sign and of $\langle +0, +0 \rangle$: a compiler that swapped the operands
would have turned \code{sqrt([-0, r])} into the empty set under AH (Apple
clang 21 and GCC 14 to its trunk keep the order, but Clang 21 merges
\code{vmaxq\_f64(x, y)} and \code{vmaxq\_f64(y, x)} into one instruction:
it takes the intrinsic for commutative). It now clears the lanes whose sign
bit is set, a negative number or $-0$, with a mask, as the part of the
interval in $\itv{0}{+\infty}$:\newinvfive{} the same bits as the maximum
with AH clear, on 20 million pairs of bounds that are no NaN and every pair
of the special values, and no dependence on AH or on the order of the
operands. The Linux and XNU kernels and the glibc, whose sources were read,
do not set AH by default; a program that does (an emulator) is outside what
GAOL supports. Modelled in software on 25.8 million results, the bounds are
those of the kernels, but for the sign of a zero.
\file{tests/rounding\_direction.cpp} runs the operations with AH set where
the processor retains the bit, and says so where it does not (the M2 and the
emulators QEMU 7.2 and 8.2 do not; QEMU 10.0 does, with
\cmd{qemu-aarch64 -cpu max}, where the test runs on 441 pairs of operands
per operation, and fails 189 checks with the operands of the former maximum
exchanged, the square root of \code{[-0, r]} being then a half-NaN
interval).
\item[$\bullet$] \textbf{The modes that make a result approximate or flush
it} are cleared by GAOL where it has a path that reads the control register:
FZ and FIZ on 64-bit ARM (FZ alone on 32-bit ARM), FTZ and DAZ on x86 (the
AVX-512 path does not read MXCSR, and clears them where a test shows them
set; \pxref{sec:ftz}). \code{NI} of FPSCR on POWER (non-IEEE mode, where
results ``may be approximate'') and \code{FS} of the FCSR on MIPS (flush
denormalized results to zero) are states of the same kind that GAOL does not
clear on the FPU intervals of these processors; a path for them would have
to.
\item[$\bullet$] \textbf{The pragma \code{FENV\_ROUND}, the
\code{\_mm256\_add\_round\_pd} of Clang 20 and the 256-bit
\code{\{ru-sae\}} assembly} compile and are wrong
(Section~\ref{sec:compilers-give}).
\item[$\bullet$] \textbf{The NEON quotient on a Neoverse N2} (the arm64
runners of GitHub, Azure Cobalt 100; the A/B of \file{performance.yml} of
7 November 2026, the medians of three runs, the base and the head giving the
same figures): \code{x / y} takes
3.96 ns with Clang, 4.04 with clang-cl, 4.49 with GCC and 4.92 with Visual
C++, where the FPU intervals take 4.29 with GCC, so that the one operation
the NEON intervals do not gain there is the quotient (with GCC \code{x * y}
takes 2.95 ns against 6.37 and \code{x + y} 1.26 against 1.87). It took 5.36
ns with clang-cl and 5.40 with Visual C++ before \code{divide\_upward()}
told the signs of the divisor on general registers (commit c42f95e, 5.35 to
3.96 ns with Clang 18), where the FPU intervals took 4.35 with clang-cl. It
was not the vector division: two scalar divisions in its place
(\code{vdiv\_f64} on each lane, the same bounds) took 6.17 and 6.30 ns, and
the time was in the rest of the kernel, about 150 instructions in
\code{div\_bounds()} of which some 55 were vector ones (counted in the
object of GCC 12.2, the common path being a part of them). On the Apple M2
(7 November 2026, mains power, three runs of each build alternated)
\code{x / y} takes 2.07 ns with the NEON intervals against 3.20 with the FPU
ones, the line of 7 quotients 18.4 ns against 32.4, and \code{x * y}
1.52 ns: the N2 takes 1.9 times as long for the product, a ratio that the
two vector pipes of the N2 against the four of the M2 would explain (not
measured: \cmd{perf} is not allowed on the runners and the machine).
\end{itemize}

\subsection{What was not measured}

No measurement of this chapter on x86 (the SSE2 and AVX-512 paths) was made on another processor than an Intel i7-1185G7, for
the times of Section~\ref{sec:computing-performance} and the divisions of
Table~\ref{tab:computing-verdicts}; no path was timed on RISC-V, POWER, s390x,
LoongArch, MIPS or WebAssembly (the tests of GAOL run on POWER, s390x and
RISC-V under QEMU in the continuous integration, \pxref{sec:ci}): the figures
given for these come from uops.info or the references, and the compilations made to study their instruction sets were made on
Compiler Explorer, without running the code. The calls
were counted on the NEON build only. GCC never ran a program of this study: the behaviour of the
intrinsics and of the pragma under it is read from the assembly. The profile
has one machine and one compiler, and \cmd{sample} is the macOS profiler.
The cost of writing FPSCR (32-bit ARM, POWER), FPC (s390x), FCSR and frm is
not in the references read, nor is the manual of MIPS. Whether FPCR.AH
changes the kernels was not measured on hardware: the M2 has no FEAT\_AFP (the
test runs with AH set under QEMU 10.0).

%%--------------------------------------------------------------------------
\section{Discussion}
%%--------------------------------------------------------------------------

This chapter has shown how GAOL v5 computes an interval: as one vector of two
doubles, \code{<-left, right>}, on which the SSE2 and the NEON kernels
compute both bounds with one or a few instructions rounded upward, with no
branch on the signs, the bounds staying in registers from one operation to
the next. What costs is not the arithmetic but the rounding direction: a
write of the control register takes about 3 ns on x86 and 4 to 11 ns on an
Apple M2, against one or two nanoseconds for the operation, and the choice
of GAOL v5 is to check the direction with one comparison and to write the
register only when the direction is not upward already. The measurements
bear the choice out. GAOL v5 is faster than GAOL 4.2.3 on every line measured and on every
operation but the square root, which takes the same time, 1.06 to 4.4 times
on x86 and 1.7 to 8.3 times on
the M2, and it is the fastest of the six libraries compared on the
arithmetic, the integer power and most formulas, with the bounds of
libieeep1788, which computes every bound with MPFR, 17 to 135 times faster.
Where it is slower, on \code{sin} and \code{cos} against filib++ and Solaris Studio, and on the square root, \code{exp}, \code{log} and the real power against PROFIL/BIAS, whose bounds are among the widest, it pays for a correct rounding over the whole range that the others
do not give: the test of enclosure found no result of GAOL v5 that misses
the value of the function, where 20 to 34\,\% of Boost.Interval's do.

\medskip
The trade-offs are those of the rounding direction. Preserving it costs
no write of the control register on the AVX-512 path, whose instructions carry their own direction,
and that is why a build takes that path with \option{GAOL\_PRESERVE\_ROUNDING}
only: without the preservation the operations take 20\,\% longer with the path at the median, and the
rule relies on one processor, with halved AVX-512 units. On 64-bit ARM, no
instruction carries its direction, and preserving it for a program that
computes to nearest makes a line of operations 13.0 to 43.5 times slower on
the M2. The paths that remain to try are few: the RISC-V scalar
instructions are the only other ones with a rounding direction by
instruction, and they need a board; POWER and s390x need hardware to
measure; the 128 and 256-bit forms of AVX-512 and the vectors of SVE offer
no embedded rounding. The cheap gains left open are a vector square root on
SSE2, a two-lane \code{mid()} and the argument reduction of the
trigonometric functions. Whether FPCR.AH changes the NEON kernels could not
be measured on hardware, the M2 having no FEAT\_AFP; the square root was made independent of it in software, and the other kernels differ under it only in the sign of a zero bound.

\medskip
These measurements rest on the tests, which check that every path gives the
same bounds in the four rounding directions of the calling code, with the
subnormals, the infinities and the empty set, and that a change does not
move the bounds by a double. The next chapter describes what the floating-point
numbers and the errors of GAOL look like to a program; the principles
and the programs of the tests, and the continuous integration that runs them
on every platform, are described in Chapter~\ref{chap:tests}.
